% Complete reviewed Unit 10; editorial master: chapters/ch10/answer-key.md
% Markdown SHA-256: f87e5d18fcabcc4502d7b2b72312c835942362e38bdb41e8a42c05b0acd29c36
% One-use initial assembly followed by complete content and page review.
\newcommand{\runningtitle}{Unit 10. Answer key}
% Layout inherited from the reviewed previous course; set a running title per unit.
% This is a small hand-maintained TeX layout, not a conversion pipeline.
\RequirePackage{fix-cm}
\documentclass[11pt,twoside]{article}
\usepackage{fontspec}
\setmainfont{Cambria}[
  Renderer=HarfBuzz,
  BoldFont=cambriab.ttf,
  ItalicFont=cambriai.ttf,
  BoldItalicFont=cambriaz.ttf
]
\setmonofont{lmmono10-regular.otf}[BoldFont=lmmonolt10-bold.otf]
\usepackage[paperwidth=165mm,paperheight=241mm,inner=21mm,outer=19mm,
  top=22mm,bottom=20mm,headheight=14pt,headsep=7mm,footskip=10mm]{geometry}
\usepackage{xcolor,array,booktabs,longtable,tabularx}
\usepackage{enumitem,fancyhdr}
\usepackage[unicode,hidelinks]{hyperref}
\definecolor{ink}{HTML}{000000}
\definecolor{muted}{HTML}{4C5558}
\setlength{\parindent}{0pt}
\setlength{\parskip}{3.5pt}
\AtBeginDocument{\fontsize{11}{14.3}\selectfont}
\setlength{\tabcolsep}{4pt}
\renewcommand{\arraystretch}{1.10}
\setlist[itemize]{leftmargin=1.3em,label=\textbullet,itemsep=2pt,topsep=3pt,parsep=0pt}
\setlist[enumerate]{leftmargin=1.8em,itemsep=3pt,topsep=4pt,parsep=0pt}
\widowpenalty=10000
\clubpenalty=10000
\emergencystretch=1em
\raggedbottom
\pagestyle{fancy}
\fancyhf{}
\providecommand{\runningtitle}{Old English}
\fancyhead[LE,RO]{\small\itshape \runningtitle}
\fancyhead[LO,RE]{\small Old English}
\fancyfoot[LE,RO]{\small\thepage}
\renewcommand{\headrulewidth}{0.3pt}
\newcommand{\booktitle}[2]{
  {\fontsize{9}{11}\selectfont\color{muted}#1\par}
  \vspace{3pt}
  {\fontsize{20}{24}\selectfont\color{ink}#2\par}
  \vspace{6pt}
}
\newcommand{\source}[1]{\par{\footnotesize\color{muted}#1\par}}
\newcommand{\readingtext}[1]{
  \par\addvspace{5pt}
  \begingroup\leftskip=1.5em
    \fontsize{12}{16}\selectfont #1\par
  \endgroup\addvspace{5pt}
}
\makeatletter
\renewcommand{\section}{\@startsection{section}{1}{\z@}
  {-2.4ex \@plus -0.6ex \@minus -0.2ex}{0.8ex \@plus 0.2ex}
  {\normalfont\fontsize{12}{15}\selectfont\bfseries\color{ink}}}
\renewcommand{\subsection}{\@startsection{subsection}{2}{\z@}
  {-1.8ex \@plus -0.5ex \@minus -0.2ex}{0.5ex \@plus 0.2ex}
  {\normalfont\fontsize{11}{14}\selectfont\bfseries\color{ink}}}
\makeatother

\hyphenpenalty=10000
\exhyphenpenalty=10000
\emergencystretch=3em
\newcommand{\itempara}[2]{\par\addvspace{3pt}\noindent\hangindent=1.8em\hangafter=1\makebox[1.8em][l]{#1.}#2\par}
\newcommand{\reserve}[1]{\par\penalty-100\begingroup\dimen0=\pagegoal\advance\dimen0 by-\pagetotal\ifdim\dimen0<#1\newpage\fi\endgroup}
\begin{document}
\addfontfeatures{Ligatures=TeXOff}
\raggedright
\booktitle{UNIT 10}{Answer key}

% mdparagraph:P001:begin
The models preserve the meaning and controls of each actual task. Documented forms and grammatical word orders can differ from a model. Where English leaves a person's gender or an addressee's number open, accept a consistent defensible choice. Distinguish valid Old English from compliance with an expressly requested construction: a sentence with \textbf{mid} may be grammatical while failing a request for bare instrumental means.
% mdparagraph:P001:end

% mdparagraph:P002:begin
Cook's English prompts below retain their original wording and identifiers. The Old English models and explanations are new teaching material. The fourteen teaching examples are not separate graded tasks; this key covers the 89 assigned exercises. [1]
% mdparagraph:P002:end

% mdparagraph:P003:begin
For future reference, a present indicative is often sufficient in context. \textbf{Willan} contributes the taught desire, willingness or refusal; it is not a mechanical replacement for English \emph{will}. Specific defensible volitional alternatives are noted below. Ordinary predicates agree with their own subjects; object descriptions agree with their objects; an internal genitive does not change the external case of its head phrase.
% mdparagraph:P003:end

% mdparagraph:P004:begin
The give-to-drink and \textbf{swā lange swā} expressions, and the Judgment reading's local helps, remain available when answering. They are not tests of unaided recall of a larger infinitive or subordinate-clause system. In complete form tasks, every labeled cell is required even where several cells have the same spelling.
% mdparagraph:P004:end

\clearpage
% mdparagraph:P005:begin
\section*{F. Complete forms}
% mdparagraph:P005:end

\reserve{9\baselineskip}
\par\begin{minipage}{\linewidth}\raggedright
% mdparagraph:P006:begin
\textbf{F01. Complete fæder.}
% mdparagraph:P006:end

% item: NEW-U10-F01
% key: KEY-NEW-U10-F01
% mdtable:T01:begin
\par\begingroup
\fontsize{10.5}{13.2}\selectfont
\setlength{\LTleft}{0pt}\setlength{\LTright}{\fill}
\noindent
\begin{tabular}{@{}>{\raggedright\arraybackslash}p{39.46mm}>{\raggedright\arraybackslash}p{39.46mm}>{\raggedright\arraybackslash}p{39.46mm}@{}}
\toprule
\textbf{Case} & \textbf{Singular} & \textbf{Plural} \\
\midrule
Nominative & fæder & fæderas \\
Genitive & fæder & fædera \\
Dative & fæder & fæderum \\
Accusative & fæder & fæderas \\
\bottomrule
\end{tabular}
\par\endgroup

% mdtable:T01:end
% mdparagraph:P007:begin
\textbf{Target and interpretation:} All four selected singular cells are endingless \textbf{fæder}. Plural nominative/\allowbreak{}accusative \textbf{fæderas}, genitive \textbf{fædera} and dative \textbf{fæderum} keep the stated masculine pattern. Documented genitive singular \textbf{fæderes} and dative \textbf{fædere} are valid variants, but should be distinguished from the selected endingless cells requested here. Do not impose the umlaut pattern of another kinship noun or invent a changed singular vowel. [2]
% mdparagraph:P007:end

\end{minipage}\par\addvspace{4pt}

\par\begin{minipage}{\linewidth}\raggedright
% mdparagraph:P008:begin
\textbf{F02. Complete gāt.}
% mdparagraph:P008:end

% item: NEW-U10-F02
% key: KEY-NEW-U10-F02
% mdtable:T02:begin
\par\begingroup
\fontsize{10.5}{13.2}\selectfont
\setlength{\LTleft}{0pt}\setlength{\LTright}{\fill}
\noindent
\begin{tabular}{@{}>{\raggedright\arraybackslash}p{39.46mm}>{\raggedright\arraybackslash}p{39.46mm}>{\raggedright\arraybackslash}p{39.46mm}@{}}
\toprule
\textbf{Case} & \textbf{Singular} & \textbf{Plural} \\
\midrule
Nominative & gāt & gǣt \\
Genitive & gāte & gāta \\
Dative & gǣt & gātum \\
Accusative & gāt & gǣt \\
\bottomrule
\end{tabular}
\par\endgroup

% mdtable:T02:end
% mdparagraph:P009:begin
\textbf{Target and interpretation:} In this selected table \textbf{gǣt} represents dative singular and nominative/\allowbreak{}accusative plural. Genitive singular is selected \textbf{gāte}. The documented alternative genitive \textbf{gǣt} and dative \textbf{gāte} are valid language; the phrase \emph{in the selected paradigm} prevents that genitive variant from becoming a fourth required answer. Keep \textbf{gāta/\allowbreak{}gātum} without umlaut, and do not invent a nominative plural in \textbf{-as}. [2]
% mdparagraph:P009:end

\end{minipage}\par\addvspace{4pt}

\par\begin{minipage}{\linewidth}\raggedright
% mdparagraph:P010:begin
\textbf{F03. Forwiernan: complete indicative.}
% mdparagraph:P010:end

% item: NEW-U10-F03
% key: KEY-NEW-U10-F03
% mdtable:T03:begin
% course-table-style: 2026-10
% course-panel:T03-P01:begin
\par\begingroup
\fontsize{10.5}{13.2}\selectfont
\setlength{\LTleft}{0pt}\setlength{\LTright}{\fill}
\noindent
\begin{tabular}{@{}>{\raggedright\arraybackslash}p{39.459mm}>{\raggedright\arraybackslash}p{39.459mm}>{\raggedright\arraybackslash}p{39.459mm}@{}}
\toprule
\textbf{Person and number} & \textbf{Present} & \textbf{Preterite} \\
\midrule
1 sg. & forwierne & forwiernde \\
2 sg. & forwiernst & forwierndest \\
3 sg. & forwiernð & forwiernde \\
1–3 pl. & forwiernað & forwierndon \\
\bottomrule
\end{tabular}
\par\endgroup
% course-panel:T03-P01:end
% mdtable:T03:end
% mdparagraph:P011:begin
\textbf{Target and interpretation:} Give all twelve labeled cells, including the three identical forms in each plural. The present uses \textbf{forwiern-}; the preterite uses dental \textbf{forwiernd-}, with second-singular \textbf{-est} and plural \textbf{-on}. For example, \textbf{forwiernst} means you deny, whereas \textbf{forwierndest} means you denied. Do not remove the past dental from a preterite or replace it with a strong-verb vowel change. Other evidenced spellings may be grammatical, but the requested table is this selected \textbf{forwiernan} system.
% mdparagraph:P011:end

\end{minipage}\par\addvspace{4pt}

\reserve{8\baselineskip}
% mdparagraph:P012:begin
\textbf{F04. Unnan: supported indicative and subjunctive.}
% mdparagraph:P012:end

% item: NEW-U10-F04
% key: KEY-NEW-U10-F04
\reserve{11\baselineskip}
% mdparagraph:P013:begin
\textbf{Indicative}
% mdparagraph:P013:end

% mdtable:T04:begin
% course-table-style: 2026-10
% course-panel:T04-P01:begin
\par\begingroup
\fontsize{10.5}{13.2}\selectfont
\setlength{\LTleft}{0pt}\setlength{\LTright}{\fill}
\noindent
\begin{tabular}{@{}>{\raggedright\arraybackslash}p{39.459mm}>{\raggedright\arraybackslash}p{39.459mm}>{\raggedright\arraybackslash}p{39.459mm}@{}}
\toprule
\textbf{Person and number} & \textbf{Present} & \textbf{Preterite} \\
\midrule
1 sg. & ann & ūðe \\
2 sg. & — & ūðest \\
3 sg. & ann & ūðe \\
1–3 pl. & unnon & ūðon \\
\bottomrule
\end{tabular}
\par\endgroup
% course-panel:T04-P01:end
% mdtable:T04:end
\reserve{11\baselineskip}
% mdparagraph:P014:begin
\textbf{Subjunctive}
% mdparagraph:P014:end

% mdtable:T05:begin
% course-table-style: 2026-10
% course-panel:T05-P01:begin
\par\begingroup
\fontsize{10.5}{13.2}\selectfont
\setlength{\LTleft}{0pt}\setlength{\LTright}{\fill}
\noindent
\begin{tabular}{@{}>{\raggedright\arraybackslash}p{39.459mm}>{\raggedright\arraybackslash}p{39.459mm}>{\raggedright\arraybackslash}p{39.459mm}@{}}
\toprule
\textbf{Person and number} & \textbf{Present} & \textbf{Preterite} \\
\midrule
1 sg. & unne & ūðe \\
2 sg. & unne & ūðe \\
3 sg. & unne & ūðe \\
1–3 pl. & unnen & ūðen \\
\bottomrule
\end{tabular}
\par\endgroup
% course-panel:T05-P01:end
% mdtable:T05:end
% mdparagraph:P015:begin
\textbf{Target and interpretation:} Present \textbf{ann/\allowbreak{}unnon} has present granting meaning; past indicative \textbf{ūðe/\allowbreak{}ūðest/\allowbreak{}ūðon} keeps its distinct endings. In the subjunctive all singular persons share \textbf{unne} or \textbf{ūðe}, and all plural persons \textbf{unnen} or \textbf{ūðen}. Present \textbf{unne} and \textbf{unnen} serve grants, wishes or addressed commands; an instruction to one addressee selects second-singular \textbf{unne}.
% mdparagraph:P015:end

% mdparagraph:P016:begin
The dash marks a present indicative second-singular cell not established by the checked sources. It is not a lexical form and does not prove that no historical variant could ever exist. Do not fill it by analogy or manufacture plural \textbf{unnað}. Cook §129 labels \textbf{unne} imperative; Fulk §123 describes subjunctive substitution where a formal imperative is absent. The course uses the supported subjunctive analysis while preserving Cook's terminology. Past subjunctive second-singular \textbf{ūðe}, unlike indicative \textbf{ūðest}, has no \textbf{-st}. [3]
% mdparagraph:P016:end

\reserve{12\baselineskip}
\par\begin{minipage}{\linewidth}\raggedright
% mdparagraph:P017:begin
\textbf{F05. Diagnostic fēdan forms.}
% mdparagraph:P017:end

% item: NEW-U10-F05
% key: KEY-NEW-U10-F05
% mdtable:T06:begin
\par\begingroup
\fontsize{10.5}{13.2}\selectfont
\setlength{\LTleft}{0pt}\setlength{\LTright}{\fill}
\noindent
\begin{tabular}{@{}>{\raggedright\arraybackslash}p{81.75mm}>{\raggedright\arraybackslash}p{39.43mm}@{}}
\toprule
\textbf{Requested cell} & \textbf{Selected form} \\
\midrule
Present indicative second singular & fētst \\
Present indicative third singular & fēt \\
Present indicative third plural & fēdað \\
Preterite indicative first singular & fēdde \\
Preterite indicative third plural & fēddon \\
Command to one addressee & fēd \\
Command to more than one & fēdað \\
\bottomrule
\end{tabular}
\par\endgroup

% mdtable:T06:end
% mdparagraph:P018:begin
\textbf{Target and interpretation:} The present second/\allowbreak{}third singular has the checked dental contraction; it is not obtained by blindly retaining \textbf{d} before an ending. The past uses \textbf{fēdd-}. The plural command and plural present indicative share \textbf{fēdað}, but their functions differ. The attested third-singular variants \textbf{fētt/\allowbreak{}fēdeð} are defensible language, while \textbf{fēt} is the selected target. Do not substitute prefixed \textbf{āfēdan} or treat singular command \textbf{fēd} as an indicative form.
% mdparagraph:P018:end

\end{minipage}\par\addvspace{4pt}

\par\begin{minipage}{\linewidth}\raggedright
% mdparagraph:P019:begin
\textbf{F06. Distinct instrumentals and dative-shaped forms.}
% mdparagraph:P019:end

% item: NEW-U10-F06
% key: KEY-NEW-U10-F06
% mdtable:T07:begin
\par\begingroup
\fontsize{10.5}{13.2}\selectfont
\setlength{\LTleft}{0pt}\setlength{\LTright}{\fill}
\noindent
\begin{tabular}{@{}>{\raggedright\arraybackslash}p{20.49mm}>{\raggedright\arraybackslash}p{23.44mm}>{\raggedright\arraybackslash}p{26.12mm}>{\raggedright\arraybackslash}p{25.5mm}>{\raggedright\arraybackslash}p{17.2mm}@{}}
\toprule
\textbf{Gender and number} & \textbf{Instrumental se} & \textbf{Strong instrumental gōd} & \textbf{Corresponding dative se} & \textbf{Strong dative gōd} \\
\midrule
Masculine singular & ðȳ; also ðon & gōde & ðǣm & gōdum \\
Neuter singular & ðȳ; also ðon & gōde & ðǣm & gōdum \\
\bottomrule
\end{tabular}
\par\endgroup

% mdtable:T07:end
% mdparagraph:P020:begin
\textbf{Target and interpretation:} Feminine singular and plural express the instrumental relation with dative-shaped forms: feminine singular \textbf{ðǣre/\allowbreak{}gōdre}, plural \textbf{ðǣm/\allowbreak{}gōdum} in all genders. Nouns likewise have no separate new ending here; the food noun \textbf{męte} shares the shape used for dative and instrumental. The construction establishes the means relation. Distinct \textbf{gōde} must not be relabeled dative \textbf{gōdum} merely because the noun's ending is shared. Bare means is deliberately practiced in this unit; \textbf{mid + dative} is the ordinary prose alternative and is grammatical where a task does not expressly demand the bare form. [4]
% mdparagraph:P020:end

\end{minipage}\par\addvspace{4pt}

\reserve{7\baselineskip}
% mdparagraph:P021:begin
\section*{A. Case and construction analysis}
% mdparagraph:P021:end

\reserve{5\baselineskip}
\par\begin{minipage}{\linewidth}\raggedright
% mdparagraph:P022:begin
\textbf{A01. Dependent and governed genitives.}
% mdparagraph:P022:end

% item: NEW-U10-A01
% key: KEY-NEW-U10-A01
% mdparagraph:P023:begin
\textbf{(a) Ðæs cuman sunu cōm.}
% mdparagraph:P023:end

% mdparagraph:P024:begin
\textbf{(b) Ic geswāc ðæs weorces.}
% mdparagraph:P024:end

% mdparagraph:P025:begin
The stranger’s son came. I ceased from the work. In the first, ðæs cuman is a genitive dependent of sunu; in the second, ðæs weorces is directly governed by geswāc.
% mdparagraph:P025:end

% mdparagraph:P026:begin
\textbf{Target and interpretation:} Relational genitive versus a lexically governed genitive. Stopped/\allowbreak{}ceased working conveys the second meaning, provided the analysis still identifies the genitive government. A genitive is not necessarily a possessor, and its case does not depend on whether English uses of.
% mdparagraph:P026:end

\end{minipage}\par\addvspace{4pt}

\par\begin{minipage}{\linewidth}\raggedright
% mdparagraph:P027:begin
\textbf{A02. A recipient with a possessor.}
% mdparagraph:P027:end

% item: NEW-U10-A02
% key: KEY-NEW-U10-A02
% mdparagraph:P028:begin
\textbf{Se munuc sealde ðæs cuman suna męte.}
% mdparagraph:P028:end

% mdparagraph:P029:begin
The monk gave food to the stranger’s son. Suna is dative singular, the recipient; męte is accusative singular, the thing given; ðæs cuman is genitive singular, the stranger whose son is meant.
% mdparagraph:P029:end

% mdparagraph:P030:begin
\textbf{Target and interpretation:} External dative head with an internal genitive dependent. Gave the stranger’s son food is equivalent. Do not make cuman the direct recipient or let its genitive replace the head noun’s dative.
% mdparagraph:P030:end

\end{minipage}\par\addvspace{4pt}

\par\begin{minipage}{\linewidth}\raggedright
% mdparagraph:P031:begin
\textbf{A03. One monk and one of the monks.}
% mdparagraph:P031:end

% item: NEW-U10-A03
% key: KEY-NEW-U10-A03
% mdparagraph:P032:begin
\textbf{(a) Ān munuc cōm.}
% mdparagraph:P032:end

% mdparagraph:P033:begin
\textbf{(b) Ān ðāra munuca cōm.}
% mdparagraph:P033:end

% mdparagraph:P034:begin
One monk came. One of the monks came. In the first ān agrees with nominative singular munuc; in the second ān stands as the nominative head and ðāra munuca is the genitive plural whole.
% mdparagraph:P034:end

% mdparagraph:P035:begin
\textbf{Target and interpretation:} Agreeing attributive numeral versus substantival numeral with genitive whole. A single monk came can express the first; the second must retain membership in the specified group. The genitive plural does not make the head or finite verb plural.
% mdparagraph:P035:end

\end{minipage}\par\addvspace{4pt}

\par\begin{minipage}{\linewidth}\raggedright
% mdparagraph:P036:begin
\textbf{A04. An agreeing numeral and a quantity word.}
% mdparagraph:P036:end

% item: NEW-U10-A04
% key: KEY-NEW-U10-A04
% mdparagraph:P037:begin
\textbf{(a) Ic hæbbe twā scēap.}
% mdparagraph:P037:end

% mdparagraph:P038:begin
\textbf{(b) Ic hæbbe fela scēapa.}
% mdparagraph:P038:end

% mdparagraph:P039:begin
I have two sheep. I have many sheep. Twā agrees with the neuter accusative plural scēap; invariant fela takes the genitive plural scēapa.
% mdparagraph:P039:end

% mdparagraph:P040:begin
\textbf{Target and interpretation:} Attributive cardinal count versus invariant fela with a genitive plural. A great many sheep is defensible for fela if the grammatical contrast is preserved. Do not carry the agreeing cardinal’s accusative construction over to fela, and do not use masculine twēgen with sheep.
% mdparagraph:P040:end

\end{minipage}\par\addvspace{4pt}

\par\begin{minipage}{\linewidth}\raggedright
% mdparagraph:P041:begin
\textbf{A05. Giving and denying.}
% mdparagraph:P041:end

% item: NEW-U10-A05
% key: KEY-NEW-U10-A05
% mdparagraph:P042:begin
\textbf{(a) Hē sealde ðǣm cuman męte.}
% mdparagraph:P042:end

% mdparagraph:P043:begin
\textbf{(b) Hē forwiernde ðǣm cuman mętes.}
% mdparagraph:P043:end

% mdparagraph:P044:begin
He gave food to the stranger. He denied the stranger food. The person is dative in both. Sęllan takes an accusative thing given; forwiernan takes a genitive thing denied.
% mdparagraph:P044:end

% mdparagraph:P045:begin
\textbf{Target and interpretation:} Dative recipient/\allowbreak{}person with accusative giving versus genitive denial. Withheld food from the stranger preserves the second meaning. The English word food looks identical but the two Old English verbs govern it differently.
% mdparagraph:P045:end

\end{minipage}\par\addvspace{4pt}

\par\begin{minipage}{\linewidth}\raggedright
% mdparagraph:P046:begin
\textbf{A06. Useful to whom?}
% mdparagraph:P046:end

% item: NEW-U10-A06
% key: KEY-NEW-U10-A06
% mdparagraph:P047:begin
\textbf{Sēo gāt is ðǣm scēaphierde behēfu.}
% mdparagraph:P047:end

% mdparagraph:P048:begin
The goat is useful to the shepherd. Gāt is feminine nominative singular, the subject; scēaphierde is masculine dative singular, the person to whom it is useful; strong behēfu is feminine nominative singular agreeing with gāt.
% mdparagraph:P048:end

% mdparagraph:P049:begin
\textbf{Target and interpretation:} Dative governed by useful, subject-predicate agreement, lexical new noun. Of use to the shepherd is equivalent. Behēfu agrees with the goat, not with the dative person; masculine behēfe would fail the selected agreement target.
% mdparagraph:P049:end

\end{minipage}\par\addvspace{4pt}

\par\begin{minipage}{\linewidth}\raggedright
% mdparagraph:P050:begin
\textbf{A07. The seeming experiencer and subject.}
% mdparagraph:P050:end

% item: NEW-U10-A07
% key: KEY-NEW-U10-A07
% mdparagraph:P051:begin
\textbf{Mē ðyncð ðis gōd.}
% mdparagraph:P051:end

% mdparagraph:P052:begin
This seems good to me. Mē is the dative experiencer. Ðis is the neuter nominative singular subject and gōd the agreeing predicate; ðyncð is third singular.
% mdparagraph:P052:end

% mdparagraph:P053:begin
\textbf{Target and interpretation:} Bounded seeming frame with a dative experiencer and neuter nominative subject. This appears good to me is equivalent. The dative experiencer is not a nominative subject. This particular seeming pattern has an expressed neuter subject; it must not be described as syntactically subjectless.
% mdparagraph:P053:end

\end{minipage}\par\addvspace{4pt}

\par\begin{minipage}{\linewidth}\raggedright
% mdparagraph:P054:begin
\textbf{A08. Two expressions of hunger.}
% mdparagraph:P054:end

% item: NEW-U10-A08
% key: KEY-NEW-U10-A08
% mdparagraph:P055:begin
\textbf{(a) Hine hyngrede.}
% mdparagraph:P055:end

% mdparagraph:P056:begin
\textbf{(b) Hē wæs hungrig.}
% mdparagraph:P056:end

% mdparagraph:P057:begin
Both mean He was hungry. Hine is accusative singular in the impersonal hunger construction, which uses third-singular hyngrede without a nominative person subject. Hē is nominative singular and controls wæs in the ordinary adjective predicate; hungrig agrees with hē.
% mdparagraph:P057:end

% mdparagraph:P058:begin
\textbf{Target and interpretation:} Accusative experiencer in an impersonal hunger clause versus nominative personal adjective predicate. He felt hungry is acceptable for the first. Do not change hyngrede to agree with an accusative experiencer as though that pronoun were the nominative subject.
% mdparagraph:P058:end

\end{minipage}\par\addvspace{4pt}

\par\begin{minipage}{\linewidth}\raggedright
% mdparagraph:P059:begin
\textbf{A09. Accompaniment and means.}
% mdparagraph:P059:end

% item: NEW-U10-A09
% key: KEY-NEW-U10-A09
% mdparagraph:P060:begin
\textbf{(a) Se munuc cōm mid ðǣm cuman.}
% mdparagraph:P060:end

% mdparagraph:P061:begin
\textbf{(b) Se munuc fēdde ðone cuman mid gōdum męte.}
% mdparagraph:P061:end

% mdparagraph:P062:begin
The monk came with the stranger. The monk fed the stranger with good food. Mid expresses accompaniment in the first and means in the second; it governs dative in both.
% mdparagraph:P062:end

% mdparagraph:P063:begin
\textbf{Target and interpretation:} Accompaniment and means expressed by the same preposition with dative. By means of good food makes the second relation explicit. Identical English with and identical prepositional case do not make the two relations identical.
% mdparagraph:P063:end

\end{minipage}\par\addvspace{4pt}

\par\begin{minipage}{\linewidth}\raggedright
% mdparagraph:P064:begin
\textbf{A10. Distinct instrumental forms.}
% mdparagraph:P064:end

% item: NEW-U10-A10
% key: KEY-NEW-U10-A10
% mdparagraph:P065:begin
\textbf{(a) Se munuc fēdde ðone cuman gōde męte.}
% mdparagraph:P065:end

% mdparagraph:P066:begin
\textbf{(b) Se munuc fēdde ðone cuman ðȳ męte.}
% mdparagraph:P066:end

% mdparagraph:P067:begin
The monk fed the stranger with good food. The monk fed the stranger with that food. Gōde is the distinct masculine singular strong instrumental; ðȳ is the distinct instrumental determiner. The noun męte has the same surface used for dative/\allowbreak{}instrumental, so its ending alone does not distinguish them.
% mdparagraph:P067:end

% mdparagraph:P068:begin
\textbf{Target and interpretation:} Bare instrumental means, distinct adjective/\allowbreak{}determiner forms and noun syncretism. That food is the required demonstrative meaning in the second; an article-like translation the food is contextually possible, but the teaching example selects that. Do not call gōde a weak nominative merely because another paradigm cell has the same spelling.
% mdparagraph:P068:end

\end{minipage}\par\addvspace{4pt}

\par\begin{minipage}{\linewidth}\raggedright
% mdparagraph:P069:begin
\textbf{A11. The perceived condition of an object.}
% mdparagraph:P069:end

% item: NEW-U10-A11
% key: KEY-NEW-U10-A11
% mdparagraph:P070:begin
\textbf{Ic geseah ðone cuman untrumne.}
% mdparagraph:P070:end

% mdparagraph:P071:begin
I saw the stranger sick. Both ðone cuman and untrumne are masculine accusative singular. The adjective describes the perceived object and agrees with it.
% mdparagraph:P071:end

% mdparagraph:P072:begin
\textbf{Target and interpretation:} Perception with a directly agreeing description of the object. I saw that the stranger was sick conveys the meaning in English, provided the Old English analysis does not invent an expressed finite ðæt clause. Untrumne does not describe the nominative speaker.
% mdparagraph:P072:end

\end{minipage}\par\addvspace{4pt}

\par\begin{minipage}{\linewidth}\raggedright
% mdparagraph:P073:begin
\textbf{A12. Naming people and describing a noun.}
% mdparagraph:P073:end

% item: NEW-U10-A12
% key: KEY-NEW-U10-A12
% mdparagraph:P074:begin
\textbf{(a) Se munuc scrȳt ðā nacodan.}
% mdparagraph:P074:end

% mdparagraph:P075:begin
\textbf{(b) Se munuc scrȳt ðone nacodan cuman.}
% mdparagraph:P075:end

% mdparagraph:P076:begin
The monk clothes the naked (people). The monk clothes the naked stranger. Nacodan is weak masculine accusative plural used substantivally in the first, but weak masculine accusative singular modifying cuman in the second. The articles and expressed noun/\allowbreak{}context distinguish reference although nacodan is syncretic.
% mdparagraph:P076:end

% mdparagraph:P077:begin
\textbf{Target and interpretation:} Substantival personal adjective versus an adjective modifying an expressed singular noun. Covers/\allowbreak{}provides clothing for is defensible when the intended clothing sense is retained. The ending -an alone does not establish number; do not supply a singular person in the first or a plural stranger in the second.
% mdparagraph:P077:end

\end{minipage}\par\addvspace{4pt}

\reserve{7\baselineskip}
% mdparagraph:P078:begin
\section*{T. Controlled composition}
% mdparagraph:P078:end

\reserve{5\baselineskip}
\par\begin{minipage}{\linewidth}\raggedright
% mdparagraph:P079:begin
\textbf{T01. Denying food.}
% mdparagraph:P079:end

% item: NEW-U10-T01
% key: KEY-NEW-U10-T01
% mdparagraph:P080:begin
\emph{The monks deny the stranger food. Use forwiernan.}
% mdparagraph:P080:end

% mdparagraph:P081:begin
\textbf{Ðā munucas forwiernað ðǣm cuman mętes.}
% mdparagraph:P081:end

% mdparagraph:P082:begin
\textbf{Target and interpretation:} \textbf{Munucas} is nominative plural and controls \textbf{forwiernað}. The stranger is dative singular \textbf{ðǣm cuman}, while the food is genitive \textbf{mętes}. The monks deny food to one stranger; they do not deny the stranger's existence or deny the food to themselves. Grammatical word-order variants are acceptable, but accusative food would miss the requested denial frame.
% mdparagraph:P082:end

\end{minipage}\par\addvspace{4pt}

\par\begin{minipage}{\linewidth}\raggedright
% mdparagraph:P083:begin
\textbf{T02. Granting food.}
% mdparagraph:P083:end

% item: NEW-U10-T02
% key: KEY-NEW-U10-T02
% mdparagraph:P084:begin
\emph{I grant the stranger food. Use unnan.}
% mdparagraph:P084:end

% mdparagraph:P085:begin
\textbf{Ic ann ðǣm cuman mętes.}
% mdparagraph:P085:end

% mdparagraph:P086:begin
\textbf{Target and interpretation:} First-singular indicative \textbf{ann} agrees with \textbf{ic} and means grant. \textbf{Ðǣm cuman} is the dative recipient, \textbf{mętes} the genitive thing granted. \textbf{Unne} would instead introduce the subjunctive form used for a wish or command. A giving sentence with \textbf{sęllan} can be grammatical but does not demonstrate the expressly requested \textbf{unnan} construction.
% mdparagraph:P086:end

\end{minipage}\par\addvspace{4pt}

\par\begin{minipage}{\linewidth}\raggedright
% mdparagraph:P087:begin
\textbf{T03. Members of a group.}
% mdparagraph:P087:end

% item: NEW-U10-T03
% key: KEY-NEW-U10-T03
% mdparagraph:P088:begin
\emph{Some of the monks came with the stranger. Which of the monks came? Use sum and hwilc with a genitive whole.}
% mdparagraph:P088:end

% mdparagraph:P089:begin
\textbf{Sume ðāra munuca cōmon mid ðǣm cuman.}
% mdparagraph:P089:end

% mdparagraph:P090:begin
\textbf{Hwilc ðāra munuca cōm?}
% mdparagraph:P090:end

% mdparagraph:P091:begin
\textbf{Target and interpretation:} \textbf{Sume} is a nominative plural head selecting genitive \textbf{ðāra munuca}, and therefore takes plural \textbf{cōmon}. \textbf{Mid ðǣm cuman} gives the accompanying stranger in the dative. In the displayed question, singular \textbf{hwilc} selects the same genitive whole and takes \textbf{cōm}. English \emph{which of the monks} does not specify how many were selected: \textbf{Hwilce ðāra munuca cōmon?} is also valid. The head, not the genitive group's plurality, controls the finite verb.
% mdparagraph:P091:end

\end{minipage}\par\addvspace{4pt}

\par\begin{minipage}{\linewidth}\raggedright
% mdparagraph:P092:begin
\textbf{T04. Many sheep.}
% mdparagraph:P092:end

% item: NEW-U10-T04
% key: KEY-NEW-U10-T04
% mdparagraph:P093:begin
\emph{The shepherd has many sheep. Use fela.}
% mdparagraph:P093:end

% mdparagraph:P094:begin
\textbf{Se scēaphierde hæfð fela scēapa.}
% mdparagraph:P094:end

% mdparagraph:P095:begin
\textbf{Target and interpretation:} The shepherd is singular subject of \textbf{hæfð}. Invariant \textbf{fela} requires genitive plural \textbf{scēapa}, meaning many sheep. An agreeing expression with \textbf{manig} may be valid Old English, but it misses the explicit \textbf{fela} target. Unchanged accusative plural \textbf{scēap} would copy the ordinary direct-object construction instead of the assigned quantity frame.
% mdparagraph:P095:end

\end{minipage}\par\addvspace{4pt}

\par\begin{minipage}{\linewidth}\raggedright
% mdparagraph:P096:begin
\textbf{T05. Bare means.}
% mdparagraph:P096:end

% item: NEW-U10-T05
% key: KEY-NEW-U10-T05
% mdparagraph:P097:begin
\emph{The monks feed the strangers with good food. Express the means without a preposition.}
% mdparagraph:P097:end

% mdparagraph:P098:begin
\textbf{Ðā munucas fēdað ðā cuman gōde męte.}
% mdparagraph:P098:end

% mdparagraph:P099:begin
\textbf{Target and interpretation:} Plural \textbf{fēdað} agrees with the monks. \textbf{Ðā cuman} is the accusative plural of the people fed, although weak \textbf{cuman} also occurs in singular oblique cells. \textbf{Gōde męte} is the requested bare means phrase: the adjective has distinct masculine singular instrumental \textbf{-e}, and the noun has the dative-shaped form. Ordinary prose \textbf{mid gōdum męte} is grammatical but misses \emph{without a preposition}. The plural participants do not make the food plural.
% mdparagraph:P099:end

\end{minipage}\par\addvspace{4pt}

\par\begin{minipage}{\linewidth}\raggedright
% mdparagraph:P100:begin
\textbf{T06. That food and the sick.}
% mdparagraph:P100:end

% item: NEW-U10-T06
% key: KEY-NEW-U10-T06
% mdparagraph:P101:begin
\emph{With that food I fed the sick. Use the distinct instrumental of se and an adjective naming the sick people.}
% mdparagraph:P101:end

% mdparagraph:P102:begin
\textbf{Ðȳ męte fēdde ic ðā untruman.}
% mdparagraph:P102:end

% mdparagraph:P103:begin
\textbf{Target and interpretation:} \textbf{Ðȳ męte} expresses the means with a distinct masculine singular instrumental determiner. \textbf{Ðā untruman} names sick people and is weak accusative plural, the people fed. \textbf{Fēdde} is first-singular preterite, agreeing with \textbf{ic} despite the initial means phrase. Instrumental \textbf{ðon} is an admitted variant, and other grammatical ordering is valid. A \textbf{mid} phrase would express means but would not demonstrate the requested distinct instrumental of \textbf{se}.
% mdparagraph:P103:end

\end{minipage}\par\addvspace{4pt}

\par\begin{minipage}{\linewidth}\raggedright
% mdparagraph:P104:begin
\textbf{T07. Plural hunger experiencers.}
% mdparagraph:P104:end

% item: NEW-U10-T07
% key: KEY-NEW-U10-T07
% mdparagraph:P105:begin
\emph{They were hungry. Use the impersonal verb hyngran.}
% mdparagraph:P105:end

% mdparagraph:P106:begin
\textbf{Hīe hyngrede.}
% mdparagraph:P106:end

% mdparagraph:P107:begin
\textbf{Target and interpretation:} \textbf{Hīe} is accusative plural, the experiencers of hunger. In this subjectless construction \textbf{hyngrede} remains third-singular preterite; it does not agree with the experiencers' plural number. An ordinary copular sentence with a correctly agreeing form of \textbf{hungrig} can be grammatical, but it misses the explicit impersonal-verb target. Do not reinterpret \textbf{hīe} as a nominative subject solely from its syncretic spelling.
% mdparagraph:P107:end

\end{minipage}\par\addvspace{4pt}

\par\begin{minipage}{\linewidth}\raggedright
% mdparagraph:P108:begin
\textbf{T08. Past seeming.}
% mdparagraph:P108:end

% item: NEW-U10-T08
% key: KEY-NEW-U10-T08
% mdparagraph:P109:begin
\emph{That seemed good to us. Use the chapter’s ðyncan construction.}
% mdparagraph:P109:end

% mdparagraph:P110:begin
\textbf{Ūs ðūhte ðæt gōd.}
% mdparagraph:P110:end

% mdparagraph:P111:begin
\textbf{Target and interpretation:} Dative plural \textbf{ūs} identifies those to whom it seemed good. \textbf{Ðæt} is neuter nominative singular, the stimulus and subject, with agreeing predicate \textbf{gōd}; \textbf{ðūhte} is singular preterite. Grammatical rearrangement is acceptable. This \textbf{ðæt} is not a content conjunction, and the plural experiencers do not require a plural finite verb.
% mdparagraph:P111:end

\end{minipage}\par\addvspace{4pt}

\reserve{7\baselineskip}
% mdparagraph:P112:begin
\section*{Cook, Group III, Exercise 13}
% mdparagraph:P112:end

\reserve{5\baselineskip}
\par\begin{minipage}{\linewidth}\raggedright
% mdparagraph:P113:begin
\textbf{III.13.1 (EX-III-13-01).}
% mdparagraph:P113:end

% item: EX-III-13-01
% key: KEY-EX-III-13-01
% mdparagraph:P114:begin
\emph{Christ is the Son of God [God’s Son].}
% mdparagraph:P114:end

% mdparagraph:P115:begin
\textbf{Crīst is Godes sunu.}
% mdparagraph:P115:end

% mdparagraph:P116:begin
\textbf{Target and interpretation:} Christ is identified as God's son. \textbf{Crīst} and predicate \textbf{sunu} are nominative singular; \textbf{Godes} is the dependent genitive. The bracketed English asks for this genitive construction, with the dependent before the head and no separate article before \textbf{sunu}. Making \textbf{sunu} accusative would treat an identifying predicate as a direct object.
% mdparagraph:P116:end

\end{minipage}\par\addvspace{4pt}

\par\begin{minipage}{\linewidth}\raggedright
% mdparagraph:P117:begin
\textbf{III.13.2 (EX-III-13-02).}
% mdparagraph:P117:end

% item: EX-III-13-02
% key: KEY-EX-III-13-02
% mdparagraph:P118:begin
\emph{Christ is our Lord.}
% mdparagraph:P118:end

% mdparagraph:P119:begin
\textbf{Crīst is ūre dryhten.}
% mdparagraph:P119:end

% mdparagraph:P120:begin
\textbf{Target and interpretation:} \textbf{Ūre} is the agreeing possessive, masculine nominative singular with predicate \textbf{dryhten}; it identifies the Lord as ours. Its ending follows the singular thing possessed, not the plural possessors. Do not substitute a personal-pronoun genitive merely because English uses \emph{our}: the assigned attributive possessive has its own agreement.
% mdparagraph:P120:end

\end{minipage}\par\addvspace{4pt}

\par\begin{minipage}{\linewidth}\raggedright
% mdparagraph:P121:begin
\textbf{III.13.3 (EX-III-13-03).}
% mdparagraph:P121:end

% item: EX-III-13-03
% key: KEY-EX-III-13-03
% mdparagraph:P122:begin
\emph{When saw I you (sing.) hungry?}
% mdparagraph:P122:end

% mdparagraph:P123:begin
\textbf{Hwonne geseah ic ðē hungrigne?}
% mdparagraph:P123:end

% mdparagraph:P124:begin
\textbf{Target and interpretation:} \textbf{Hwonne} asks when; \textbf{geseah} is first-singular preterite, with \textbf{ic} as the perceiver. \textbf{Ðē} is the accusative person seen, and strong \textbf{hungrigne} is an accusative masculine singular description of that person. The model assumes a male addressee. \textbf{Hungrige}, accusative feminine singular, is valid for a female addressee: the English fixes number but not gender. A nominative adjective would wrongly describe the speaker or fail to agree with the object.
% mdparagraph:P124:end

\end{minipage}\par\addvspace{4pt}

\par\begin{minipage}{\linewidth}\raggedright
% mdparagraph:P125:begin
\textbf{III.13.4 (EX-III-13-04).}
% mdparagraph:P125:end

% item: EX-III-13-04
% key: KEY-EX-III-13-04
% mdparagraph:P126:begin
\emph{Gold is precious.}
% mdparagraph:P126:end

% mdparagraph:P127:begin
\textbf{Gold is dēorwierðe.}
% mdparagraph:P127:end

% mdparagraph:P128:begin
\textbf{Target and interpretation:} Mass \textbf{gold} is neuter nominative singular, and \textbf{dēorwierðe} is its strong nominative neuter singular predicate: gold is valuable. The final \textbf{-e} belongs to this adjective's lexical pattern; it does not make the form weak. Do not remove it to imitate the endingless neuter of \textbf{gōd}.
% mdparagraph:P128:end

\end{minipage}\par\addvspace{4pt}

\par\begin{minipage}{\linewidth}\raggedright
% mdparagraph:P129:begin
\textbf{III.13.5 (EX-III-13-05).}
% mdparagraph:P129:end

% item: EX-III-13-05
% key: KEY-EX-III-13-05
% mdparagraph:P130:begin
\emph{The Day of Judgment will come.}
% mdparagraph:P130:end

% mdparagraph:P131:begin
\textbf{Dōmes dæg cymð.}
% mdparagraph:P131:end

% mdparagraph:P132:begin
\textbf{Target and interpretation:} The day, \textbf{dæg}, is the nominative subject; \textbf{dōmes} is its genitive dependent, identifying the judgment day. Present third-singular \textbf{cymð} receives future reference from the sentence's context. A future event does not require \textbf{willan}, whose taught meaning adds desire or willingness. The genitive phrase does not make the subject plural or turn \textbf{dæg} into genitive.
% mdparagraph:P132:end

\end{minipage}\par\addvspace{4pt}

\par\begin{minipage}{\linewidth}\raggedright
% mdparagraph:P133:begin
\textbf{III.13.6 (EX-III-13-06).}
% mdparagraph:P133:end

% item: EX-III-13-06
% key: KEY-EX-III-13-06
% mdparagraph:P134:begin
\emph{The farmer says to the workman, I will sell this house.}
% mdparagraph:P134:end

% mdparagraph:P135:begin
\textbf{Se ierðling sægð ðǣm wyrhtan: “ic cīepe ðis hūs.”}
% mdparagraph:P135:end

% mdparagraph:P136:begin
\textbf{Target and interpretation:} The farmer speaks to dative \textbf{ðǣm wyrhtan}. Inside the quotation, \textbf{ic} is the farmer and \textbf{cīepe} is first singular; \textbf{ðis hūs} is the neuter accusative thing to be sold. The present can state the intended later sale. If \emph{will} expresses willingness, \textbf{ic wille ðis hūs cīepan} is also defensible, using the separately taught plain infinitive. Do not retain third-person agreement inside the farmer's first-person speech or put the speaking recipient in the accusative.
% mdparagraph:P136:end

\end{minipage}\par\addvspace{4pt}

\par\begin{minipage}{\linewidth}\raggedright
% mdparagraph:P137:begin
\textbf{III.13.7 (EX-III-13-07).}
% mdparagraph:P137:end

% item: EX-III-13-07
% key: KEY-EX-III-13-07
% mdparagraph:P138:begin
\emph{The workman answers the farmer, I will buy your house.}
% mdparagraph:P138:end

% mdparagraph:P139:begin
\textbf{Se wyrhta andswarað ðǣm ierðlinge: “ic bycge ðīn hūs.”}
% mdparagraph:P139:end

% mdparagraph:P140:begin
\textbf{Target and interpretation:} \textbf{Andswarian} takes the answered person in the dative: \textbf{ðǣm ierðlinge}. The workman is \textbf{ic} in the quotation; the farmer is the singular possessor named by \textbf{ðīn}, which agrees with accusative neuter singular \textbf{hūs}. The model's present expresses the prospective purchase. \textbf{Ic wille ðīn hūs bycgan} is valid when willingness or intention is meant. Confusing these speech roles would make the workman buy his own house.
% mdparagraph:P140:end

\end{minipage}\par\addvspace{4pt}

\par\begin{minipage}{\linewidth}\raggedright
% mdparagraph:P141:begin
\textbf{III.13.8 (EX-III-13-08).}
% mdparagraph:P141:end

% item: EX-III-13-08
% key: KEY-EX-III-13-08
% mdparagraph:P142:begin
\emph{Life is more than raiment.}
% mdparagraph:P142:end

% mdparagraph:P143:begin
\textbf{Līf is māre ðonne rēaf.}
% mdparagraph:P143:end

% mdparagraph:P144:begin
\textbf{Target and interpretation:} Life is greater in worth than clothing. \textbf{Māre} is the weak comparative, neuter nominative singular with \textbf{līf}; comparatives remain weak without an article. \textbf{Rēaf} is the parallel nominative term after \textbf{ðonne}, with the copula understood. Here \emph{more} evaluates worth, not the number of lives or garments. Do not replace \textbf{māre} by a plural quantity expression.
% mdparagraph:P144:end

\end{minipage}\par\addvspace{4pt}

\par\begin{minipage}{\linewidth}\raggedright
% mdparagraph:P145:begin
\textbf{III.13.9 (EX-III-13-09).}
% mdparagraph:P145:end

% item: EX-III-13-09
% key: KEY-EX-III-13-09
% mdparagraph:P146:begin
\emph{The stranger is in prison.}
% mdparagraph:P146:end

% mdparagraph:P147:begin
\textbf{Se cuma is on cwearterne.}
% mdparagraph:P147:end

% mdparagraph:P148:begin
\textbf{Target and interpretation:} The stranger is located in prison: \textbf{se cuma} is nominative singular, and locative \textbf{on} takes dative \textbf{cwearterne}. The English institutional expression has no article, and the model needs none before the prison. \textbf{Cweartern} without its dative ending would fail this selected location construction; \textbf{cuman} would be an oblique form of the stranger rather than the subject.
% mdparagraph:P148:end

\end{minipage}\par\addvspace{4pt}

\par\begin{minipage}{\linewidth}\raggedright
% mdparagraph:P149:begin
\textbf{III.13.10 (EX-III-13-10).}
% mdparagraph:P149:end

% item: EX-III-13-10
% key: KEY-EX-III-13-10
% mdparagraph:P150:begin
\emph{The merchant is in the inn.}
% mdparagraph:P150:end

% mdparagraph:P151:begin
\textbf{Se mangere is on ðǣm giesthūse.}
% mdparagraph:P151:end

% mdparagraph:P152:begin
\textbf{Target and interpretation:} The merchant is the subject; the inn is a location, \textbf{on ðǣm giesthūse}, with neuter dative singular agreement. \textbf{Giesthūse} retains the compound's noun stem and adds the dative ending. Nominative/\allowbreak{}accusative \textbf{giesthūs} is not the required cell after locative \textbf{on}.
% mdparagraph:P152:end

\end{minipage}\par\addvspace{4pt}

\par\begin{minipage}{\linewidth}\raggedright
% mdparagraph:P153:begin
\textbf{III.13.11 (EX-III-13-11).}
% mdparagraph:P153:end

% item: EX-III-13-11
% key: KEY-EX-III-13-11
% mdparagraph:P154:begin
\emph{The father clothes the son.}
% mdparagraph:P154:end

% mdparagraph:P155:begin
\textbf{Se fæder scrȳt ðone sunu.}
% mdparagraph:P155:end

% mdparagraph:P156:begin
\textbf{Target and interpretation:} Nominative \textbf{fæder} clothes accusative \textbf{ðone sunu}. \textbf{Scrȳt} is the selected third-singular present of \textbf{scrȳdan}, not an instruction to clothe. Accusative singular \textbf{sunu} retains lexical \textbf{-u}; dative \textbf{suna} would not express the person clothed in this frame. The identical surface \textbf{fæder} in several cases does not determine its role without the clause.
% mdparagraph:P156:end

\end{minipage}\par\addvspace{4pt}

\par\begin{minipage}{\linewidth}\raggedright
% mdparagraph:P157:begin
\textbf{III.13.12 (EX-III-13-12).}
% mdparagraph:P157:end

% item: EX-III-13-12
% key: KEY-EX-III-13-12
% mdparagraph:P158:begin
\emph{This world is great.}
% mdparagraph:P158:end

% mdparagraph:P159:begin
\textbf{Ðēos woruld is miclu.}
% mdparagraph:P159:end

% mdparagraph:P160:begin
\textbf{Target and interpretation:} \textbf{Woruld} is feminine nominative singular, so \emph{this} is \textbf{ðēos} and the strong predicate \emph{great} is selected \textbf{miclu}. Both agree with the world, not with an assumed neuter English \emph{it}. \textbf{Ðis} or masculine \textbf{micel} would miss that grammatical gender; no weak predicate is required merely because the subject has a demonstrative.
% mdparagraph:P160:end

\end{minipage}\par\addvspace{4pt}

\par\begin{minipage}{\linewidth}\raggedright
% mdparagraph:P161:begin
\textbf{III.13.13 (EX-III-13-13).}
% mdparagraph:P161:end

% item: EX-III-13-13
% key: KEY-EX-III-13-13
% mdparagraph:P162:begin
\emph{The shepherd feeds the sheep (sing.).}
% mdparagraph:P162:end

% mdparagraph:P163:begin
\textbf{Se scēaphierde fēt ðæt scēap.}
% mdparagraph:P163:end

% mdparagraph:P164:begin
\textbf{Target and interpretation:} The shepherd feeds one sheep. \textbf{Fēt} is third-singular present, and \textbf{ðæt scēap} is neuter accusative singular, the being fed. The noun's unchanged spelling can also occur in plural cells, but the English explicitly requires singular here. \textbf{Ðā scēap} would make the sheep plural and miss that control.
% mdparagraph:P164:end

\end{minipage}\par\addvspace{4pt}

\par\begin{minipage}{\linewidth}\raggedright
% mdparagraph:P165:begin
\textbf{III.13.14 (EX-III-13-14).}
% mdparagraph:P165:end

% item: EX-III-13-14
% key: KEY-EX-III-13-14
% mdparagraph:P166:begin
\emph{The merchants are rich.}
% mdparagraph:P166:end

% mdparagraph:P167:begin
\textbf{Ðā mangeras sindon welige.}
% mdparagraph:P167:end

% mdparagraph:P168:begin
\textbf{Target and interpretation:} \textbf{Mangeras} is nominative masculine plural; plural \textbf{sindon} and strong predicate \textbf{welige} agree with the merchants. The article belongs to the subject phrase and does not make the separate predicate adjective weak. Singular \textbf{is} loses the number; weak \textbf{weligan} does not meet the ordinary positive-predicate construction.
% mdparagraph:P168:end

\end{minipage}\par\addvspace{4pt}

\par\begin{minipage}{\linewidth}\raggedright
% mdparagraph:P169:begin
\textbf{III.13.15 (EX-III-13-15).}
% mdparagraph:P169:end

% item: EX-III-13-15
% key: KEY-EX-III-13-15
% mdparagraph:P170:begin
\emph{The soldier drinks wine.}
% mdparagraph:P170:end

% mdparagraph:P171:begin
\textbf{Se cęmpa drincð wīn.}
% mdparagraph:P171:end

% mdparagraph:P172:begin
\textbf{Target and interpretation:} Nominative \textbf{se cęmpa} is the drinker, and singular mass \textbf{wīn} is accusative, the drink. \textbf{Drincð} is present third singular. The soldier's weak-noun nominative \textbf{-a} is not a feminine or plural marker, and a plural verb would not agree with this subject.
% mdparagraph:P172:end

\end{minipage}\par\addvspace{4pt}

\par\begin{minipage}{\linewidth}\raggedright
% mdparagraph:P173:begin
\textbf{III.13.16 (EX-III-13-16).}
% mdparagraph:P173:end

% item: EX-III-13-16
% key: KEY-EX-III-13-16
% mdparagraph:P174:begin
\emph{The monks drank water.}
% mdparagraph:P174:end

% mdparagraph:P175:begin
\textbf{Ðā munucas druncon wæter.}
% mdparagraph:P175:end

% mdparagraph:P176:begin
\textbf{Target and interpretation:} The monks drank water in the past. \textbf{Munucas} is nominative plural; \textbf{druncon} uses the plural preterite stem and ending. \textbf{Wæter} is the singular mass object. Singular \textbf{dranc} would fail the subject's number, while present \textbf{drincað} would lose the past time.
% mdparagraph:P176:end

\end{minipage}\par\addvspace{4pt}

\clearpage
% mdparagraph:P177:begin
\section*{Cook, Group III, Exercise 14}
% mdparagraph:P177:end

\reserve{5\baselineskip}
\par\begin{minipage}{\linewidth}\raggedright
% mdparagraph:P178:begin
\textbf{III.14.1 (EX-III-14-01).}
% mdparagraph:P178:end

% item: EX-III-14-01
% key: KEY-EX-III-14-01
% mdparagraph:P179:begin
\emph{The king sat on the throne.}
% mdparagraph:P179:end

% mdparagraph:P180:begin
\textbf{Se cyning sæt on ðǣm setle.}
% mdparagraph:P180:end

% mdparagraph:P181:begin
\textbf{Target and interpretation:} \textbf{Sæt} is the third-singular preterite of \textbf{sittan}. The throne is a location expressed by neuter dative \textbf{ðǣm setle} after \textbf{on}. The king is sitting, not being seated by an unnamed agent. Accusative \textbf{setl} would miss this ordinary locative frame.
% mdparagraph:P181:end

\end{minipage}\par\addvspace{4pt}

\par\begin{minipage}{\linewidth}\raggedright
% mdparagraph:P182:begin
\textbf{III.14.2 (EX-III-14-02).}
% mdparagraph:P182:end

% item: EX-III-14-02
% key: KEY-EX-III-14-02
% mdparagraph:P183:begin
\emph{Stand on my right hand.}
% mdparagraph:P183:end

% mdparagraph:P184:begin
\textbf{Stand on mīne swīðran hand.}
% mdparagraph:P184:end

% mdparagraph:P185:begin
\textbf{Target and interpretation:} The command places its addressee at the speaker's right side. The bounded expression \textbf{on mīne swīðran hand} has feminine accusative singular \textbf{hand}, with agreeing \textbf{mīne} and weak \textbf{swīðran}. The checked source uses the same accusative side-expression with static \textbf{standað}; rest alone does not require a dative. The ending \textbf{-an} on \textbf{swīðran} by itself cannot decide the case. The English leaves the addressee's number open, so plural \textbf{Standað} is also valid. [4]
% mdparagraph:P185:end

\end{minipage}\par\addvspace{4pt}

\par\begin{minipage}{\linewidth}\raggedright
% mdparagraph:P186:begin
\textbf{III.14.3 (EX-III-14-03).}
% mdparagraph:P186:end

% item: EX-III-14-03
% key: KEY-EX-III-14-03
% mdparagraph:P187:begin
\emph{All nations shall come to the judgment.}
% mdparagraph:P187:end

% mdparagraph:P188:begin
\textbf{Ealla ðēoda cumað tō ðǣm dōme.}
% mdparagraph:P188:end

% mdparagraph:P189:begin
\textbf{Target and interpretation:} All nations are the feminine nominative plural subject, \textbf{ealla ðēoda}, with plural \textbf{cumað}. Destination \textbf{tō} governs dative \textbf{ðǣm dōme}. The context makes the present future in reference. \textbf{Ealla} is the selected feminine plural; the source also attests feminine \textbf{ealle}, so that variant need not be mistaken for masculine agreement. A singular verb or an accusative destination misses the selected frame.
% mdparagraph:P189:end

\end{minipage}\par\addvspace{4pt}

\par\begin{minipage}{\linewidth}\raggedright
% mdparagraph:P190:begin
\textbf{III.14.4 (EX-III-14-04).}
% mdparagraph:P190:end

% item: EX-III-14-04
% key: KEY-EX-III-14-04
% mdparagraph:P191:begin
\emph{The devil is (om. the) father of the unrighteous.}
% mdparagraph:P191:end

% mdparagraph:P192:begin
\textbf{Se dēofol is ðāra unrihtwīsra fæder.}
% mdparagraph:P192:end

% mdparagraph:P193:begin
\textbf{Target and interpretation:} The devil is identified as father of the unrighteous people. Substantival \textbf{ðāra unrihtwīsra} is genitive plural dependent on nominative singular \textbf{fæder}; its article belongs to the dependent people, not to the father. Follow the explicit omission of English \emph{the} before \emph{father}. An accusative plural would treat the unrighteous as an object, while an extra article on the head would miss the assigned dependent-first genitive pattern.
% mdparagraph:P193:end

\end{minipage}\par\addvspace{4pt}

\par\begin{minipage}{\linewidth}\raggedright
% mdparagraph:P194:begin
\textbf{III.14.5 (EX-III-14-05).}
% mdparagraph:P194:end

% item: EX-III-14-05
% key: KEY-EX-III-14-05
% mdparagraph:P195:begin
\emph{All men are sinful.}
% mdparagraph:P195:end

% mdparagraph:P196:begin
\textbf{Ealle męnn sindon synfulle.}
% mdparagraph:P196:end

% mdparagraph:P197:begin
\textbf{Target and interpretation:} \textbf{Ealle męnn} is the nominative masculine plural subject; \textbf{sindon} and strong predicate \textbf{synfulle} agree with it. The sentence describes all men as sinful; it does not use an adjective by itself to name \emph{the sinful}. Keep the umlauted noun plural \textbf{męnn}, and do not substitute the weak substantive form for this ordinary predicate.
% mdparagraph:P197:end

\end{minipage}\par\addvspace{4pt}

\par\begin{minipage}{\linewidth}\raggedright
% mdparagraph:P198:begin
\textbf{III.14.6 (EX-III-14-06).}
% mdparagraph:P198:end

% item: EX-III-14-06
% key: KEY-EX-III-14-06
% mdparagraph:P199:begin
\emph{The priest is sick.}
% mdparagraph:P199:end

% mdparagraph:P200:begin
\textbf{Se mæsseprēost is untrum.}
% mdparagraph:P200:end

% mdparagraph:P201:begin
\textbf{Target and interpretation:} The priest is the nominative subject, and \textbf{untrum} is his strong masculine nominative singular predicate. \textbf{Is} states present illness. This is an ordinary copular sentence: English illness does not by itself require an impersonal verb or an accusative experiencer. A visitation verb would introduce an action absent from the prompt.
% mdparagraph:P201:end

\end{minipage}\par\addvspace{4pt}

\par\begin{minipage}{\linewidth}\raggedright
% mdparagraph:P202:begin
\textbf{III.14.7 (EX-III-14-07).}
% mdparagraph:P202:end

% item: EX-III-14-07
% key: KEY-EX-III-14-07
% mdparagraph:P203:begin
\emph{Grant (sing.) me garments.}
% mdparagraph:P203:end

% mdparagraph:P204:begin
\textbf{Unne mē wǣda.}
% mdparagraph:P204:end

% mdparagraph:P205:begin
\textbf{Target and interpretation:} The command addresses one person. \textbf{Unne} is the selected present subjunctive second singular used as an addressed command; \textbf{mē} is dative recipient and \textbf{wǣda} genitive plural thing granted. Cook calls \textbf{unne} imperative; the selected analysis follows the separately evidenced subjunctive substitution. These descriptions preserve the same supported form. Do not invent a regular imperative or copy the accusative-thing government of ordinary \textbf{sęllan}. [3]
% mdparagraph:P205:end

\end{minipage}\par\addvspace{4pt}

\par\begin{minipage}{\linewidth}\raggedright
% mdparagraph:P206:begin
\textbf{III.14.8 (EX-III-14-08).}
% mdparagraph:P206:end

% item: EX-III-14-08
% key: KEY-EX-III-14-08
% mdparagraph:P207:begin
\emph{The horse is useful to the king.}
% mdparagraph:P207:end

% mdparagraph:P208:begin
\textbf{Ðæt hors is ðǣm cyninge behēfe.}
% mdparagraph:P208:end

% mdparagraph:P209:begin
\textbf{Target and interpretation:} The horse is useful; the king is the person to whom it is useful. Thus \textbf{ðæt hors} is neuter nominative singular and \textbf{ðǣm cyninge} masculine dative singular. \textbf{Behēfe} is a neuter nominative predicate agreeing with the horse, even though a masculine cell has the same spelling. The dative king is neither the subject nor an accusative object.
% mdparagraph:P209:end

\end{minipage}\par\addvspace{4pt}

\par\begin{minipage}{\linewidth}\raggedright
% mdparagraph:P210:begin
\textbf{III.14.9 (EX-III-14-09).}
% mdparagraph:P210:end

% item: EX-III-14-09
% key: KEY-EX-III-14-09
% mdparagraph:P211:begin
\emph{The loss of all your things was strange.}
% mdparagraph:P211:end

% mdparagraph:P212:begin
\textbf{Ealra ðīnra ðinga lyre wæs sellic.}
% mdparagraph:P212:end

% mdparagraph:P213:begin
\textbf{Target and interpretation:} The singular loss, \textbf{lyre}, is the subject of \textbf{wæs} and is described by strong masculine singular \textbf{sellic}. Within its genitive dependent, \textbf{ealra} and \textbf{ðīnra} agree with neuter plural \textbf{ðinga}: all the things belonged to the addressee. They do not agree with singular \textbf{lyre}. English \emph{your} leaves the possessor's number open; \textbf{ēowerra} for plural possessors, or \textbf{incerra} for two, is defensible in place of \textbf{ðīnra}. Plural things do not make the loss or its verb plural.
% mdparagraph:P213:end

\end{minipage}\par\addvspace{4pt}

\par\begin{minipage}{\linewidth}\raggedright
% mdparagraph:P214:begin
\textbf{III.14.10 (EX-III-14-10).}
% mdparagraph:P214:end

% item: EX-III-14-10
% key: KEY-EX-III-14-10
% mdparagraph:P215:begin
\emph{He departed from her.}
% mdparagraph:P215:end

% mdparagraph:P216:begin
\textbf{Hē gewāt fram hiere.}
% mdparagraph:P216:end

% mdparagraph:P217:begin
\textbf{Target and interpretation:} \textbf{Gewāt} is the singular preterite meaning departed. \textbf{Fram} governs dative \textbf{hiere}, the person from whom he departed. This \textbf{hiere} is not a possessor attached to an omitted noun. Do not analyze the prefix \textbf{ge-} as a separate perfect auxiliary or change the pronoun to an accusative merely because English says \emph{her}.
% mdparagraph:P217:end

\end{minipage}\par\addvspace{4pt}

\par\begin{minipage}{\linewidth}\raggedright
% mdparagraph:P218:begin
\textbf{III.14.11 (EX-III-14-11).}
% mdparagraph:P218:end

% item: EX-III-14-11
% key: KEY-EX-III-14-11
% mdparagraph:P219:begin
\emph{Give (sing.) me to drink.}
% mdparagraph:P219:end

% mdparagraph:P220:begin
\textbf{Sęle mē drincan.}
% mdparagraph:P220:end

% mdparagraph:P221:begin
\textbf{Target and interpretation:} Singular \textbf{sęle} gives something to the dative speaker \textbf{mē}; plain \textbf{drincan} supplies the action provided. The locally supplied idiom is \emph{give someone to drink}, with no \textbf{tō}. The infinitive remains uninflected. \textbf{Tō drincanne} can serve a purpose construction elsewhere, but substituting it here would miss the supplied idiom; unaided recall of that idiom is not being tested. [4]
% mdparagraph:P221:end

\end{minipage}\par\addvspace{4pt}

\par\begin{minipage}{\linewidth}\raggedright
% mdparagraph:P222:begin
\textbf{III.14.12 (EX-III-14-12).}
% mdparagraph:P222:end

% item: EX-III-14-12
% key: KEY-EX-III-14-12
% mdparagraph:P223:begin
\emph{Visit (plur.) the sick.}
% mdparagraph:P223:end

% mdparagraph:P224:begin
\textbf{Genēosiað ðā untruman.}
% mdparagraph:P224:end

% mdparagraph:P225:begin
\textbf{Target and interpretation:} \textbf{Genēosiað} is the required plural command, addressing more than one person. \textbf{Ðā untruman} names the sick people and is weak accusative plural, the persons visited. The checked \textbf{genēosian} frame governs an accusative person; a dative cannot be inferred from an English paraphrase such as \emph{attend to}. A singular command misses the explicit plural control, independently of the number of sick people.
% mdparagraph:P225:end

\end{minipage}\par\addvspace{4pt}

\par\begin{minipage}{\linewidth}\raggedright
% mdparagraph:P226:begin
\textbf{III.14.13 (EX-III-14-13).}
% mdparagraph:P226:end

% item: EX-III-14-13
% key: KEY-EX-III-14-13
% mdparagraph:P227:begin
\emph{When did you (sing.) see the goat?}
% mdparagraph:P227:end

% mdparagraph:P228:begin
\textbf{Hwonne gesāwe ðū ðā gāt?}
% mdparagraph:P228:end

% mdparagraph:P229:begin
\textbf{Target and interpretation:} \textbf{Gesāwe} is the second-singular preterite agreeing with \textbf{ðū}. Feminine accusative singular \textbf{ðā gāt} is the goat seen. In this object position \textbf{gǣt} would be plural, although that spelling has other roles in the full paradigm. Do not use first/\allowbreak{}third-singular \textbf{geseah} for the explicitly singular second-person subject or add English auxiliary \emph{did}.
% mdparagraph:P229:end

\end{minipage}\par\addvspace{4pt}

\par\begin{minipage}{\linewidth}\raggedright
% mdparagraph:P230:begin
\textbf{III.14.14 (EX-III-14-14).}
% mdparagraph:P230:end

% item: EX-III-14-14
% key: KEY-EX-III-14-14
% mdparagraph:P231:begin
\emph{When was he in the inn?}
% mdparagraph:P231:end

% mdparagraph:P232:begin
\textbf{Hwonne wæs hē on ðǣm giesthūse?}
% mdparagraph:P232:end

% mdparagraph:P233:begin
\textbf{Target and interpretation:} \textbf{Hwonne} asks for the time of his past presence; \textbf{wæs} agrees with singular \textbf{hē}. The inn remains dative \textbf{ðǣm giesthūse} after locative \textbf{on}. The question word neither governs the inn nor changes a location into an accusative phrase. Present \textbf{is} would change the time asked about.
% mdparagraph:P233:end

\end{minipage}\par\addvspace{4pt}

\par\begin{minipage}{\linewidth}\raggedright
% mdparagraph:P234:begin
\textbf{III.14.15 (EX-III-14-15).}
% mdparagraph:P234:end

% item: EX-III-14-15
% key: KEY-EX-III-14-15
% mdparagraph:P235:begin
\emph{This is eternal life.}
% mdparagraph:P235:end

% mdparagraph:P236:begin
\textbf{Ðis is ēce līf.}
% mdparagraph:P236:end

% mdparagraph:P237:begin
\textbf{Target and interpretation:} Neuter nominative \textbf{ðis} is identified with \textbf{ēce līf}. \textbf{Ēce} agrees with nominative neuter singular \textbf{līf} and retains its lexical \textbf{-e} in the strong form. Do not infer weak inflection from that final vowel or translate the demonstrative as a clause-introducing \emph{that}.
% mdparagraph:P237:end

\end{minipage}\par\addvspace{4pt}

\par\begin{minipage}{\linewidth}\raggedright
% mdparagraph:P238:begin
\textbf{III.14.16 (EX-III-14-16).}
% mdparagraph:P238:end

% item: EX-III-14-16
% key: KEY-EX-III-14-16
% mdparagraph:P239:begin
\emph{The prisons were empty.}
% mdparagraph:P239:end

% mdparagraph:P240:begin
\textbf{Ðā cweartern wǣron ǣmtig.}
% mdparagraph:P240:end

% mdparagraph:P241:begin
\textbf{Target and interpretation:} Endingless \textbf{cweartern} is nominative neuter plural here, established by \textbf{ðā}, plural \textbf{wǣron}, and the English. Selected strong neuter nominative plural \textbf{ǣmtig} is also endingless. Neither noun nor adjective must acquire \textbf{-as} or \textbf{-e} simply to look plural. Singular \textbf{wæs} would lose the stated number.
% mdparagraph:P241:end

\end{minipage}\par\addvspace{4pt}

\par\begin{minipage}{\linewidth}\raggedright
% mdparagraph:P242:begin
\textbf{III.14.17 (EX-III-14-17).}
% mdparagraph:P242:end

% item: EX-III-14-17
% key: KEY-EX-III-14-17
% mdparagraph:P243:begin
\emph{The living God is eternal.}
% mdparagraph:P243:end

% mdparagraph:P244:begin
\textbf{Se cwica God is ēce.}
% mdparagraph:P244:end

% mdparagraph:P245:begin
\textbf{Target and interpretation:} \textbf{Cwica} is weak masculine nominative singular inside the definite phrase \textbf{se cwica God}. \textbf{Ēce} is a separate strong predicate describing God as eternal. The two adjectives have different functions: the subject's article controls the ordinary attributive adjective, not the independent predicate. Using strong \textbf{cwic} inside this selected definite phrase would miss the requested agreement pattern.
% mdparagraph:P245:end

\end{minipage}\par\addvspace{4pt}

\reserve{7\baselineskip}
% mdparagraph:P246:begin
\section*{Cook, Group III, Exercise 15}
% mdparagraph:P246:end

\reserve{5\baselineskip}
\par\begin{minipage}{\linewidth}\raggedright
% mdparagraph:P247:begin
\textbf{III.15.1 (EX-III-15-01).}
% mdparagraph:P247:end

% item: EX-III-15-01
% key: KEY-EX-III-15-01
% mdparagraph:P248:begin
\emph{He stood before the king.}
% mdparagraph:P248:end

% mdparagraph:P249:begin
\textbf{Hē stōd beforan ðǣm cyninge.}
% mdparagraph:P249:end

% mdparagraph:P250:begin
\textbf{Target and interpretation:} \textbf{Stōd} is singular preterite: he stood in front of the king. The selected locative preposition \textbf{beforan} takes masculine dative \textbf{ðǣm cyninge}. The spelling is the joined course form. An accusative king would miss this specific frame; the case is learned with the preposition, not inferred from the modern word \emph{before} alone.
% mdparagraph:P250:end

\end{minipage}\par\addvspace{4pt}

\par\begin{minipage}{\linewidth}\raggedright
% mdparagraph:P251:begin
\textbf{III.15.2 (EX-III-15-02).}
% mdparagraph:P251:end

% item: EX-III-15-02
% key: KEY-EX-III-15-02
% mdparagraph:P252:begin
\emph{The goat eats grass.}
% mdparagraph:P252:end

% mdparagraph:P253:begin
\textbf{Sēo gāt itt gærs.}
% mdparagraph:P253:end

% mdparagraph:P254:begin
\textbf{Target and interpretation:} \textbf{Sēo gāt} is feminine nominative singular, and \textbf{itt} is the selected third-singular present of \textbf{etan}. Mass \textbf{gærs} is the accusative thing eaten. Do not confuse present \textbf{itt} with preterite \textbf{ǣt}, or treat the feminine goat as masculine because English ordinarily uses \emph{it} for an animal.
% mdparagraph:P254:end

\end{minipage}\par\addvspace{4pt}

\par\begin{minipage}{\linewidth}\raggedright
% mdparagraph:P255:begin
\textbf{III.15.3 (EX-III-15-03).}
% mdparagraph:P255:end

% item: EX-III-15-03
% key: KEY-EX-III-15-03
% mdparagraph:P256:begin
\emph{The angel flies to the earth.}
% mdparagraph:P256:end

% mdparagraph:P257:begin
\textbf{Se ęngel flīehð tō ðǣre eorðan.}
% mdparagraph:P257:end

% mdparagraph:P258:begin
\textbf{Target and interpretation:} The angel is the singular subject of selected \textbf{flīehð}; \textbf{tō} gives the destination with dative \textbf{ðǣre eorðan}. \textbf{Ęngel} means angel, distinct from the earlier fish-hook word \textbf{angel}. Preserve the checked present stem and the destination's dative; replacing the latter with an accusative imports an untaught government rule.
% mdparagraph:P258:end

\end{minipage}\par\addvspace{4pt}

\par\begin{minipage}{\linewidth}\raggedright
% mdparagraph:P259:begin
\textbf{III.15.4 (EX-III-15-04).}
% mdparagraph:P259:end

% item: EX-III-15-04
% key: KEY-EX-III-15-04
% mdparagraph:P260:begin
\emph{You (sing.) deny me food.}
% mdparagraph:P260:end

% mdparagraph:P261:begin
\textbf{Ðū forwiernst mē mętes.}
% mdparagraph:P261:end

% mdparagraph:P262:begin
\textbf{Target and interpretation:} \textbf{Forwiernst} agrees with singular \textbf{ðū}. \textbf{Mē} is the dative person denied, and \textbf{mętes} the genitive thing denied. English puts both after \emph{deny}, but their Old English cases differ. Accusative \textbf{męte}, as in ordinary giving, would miss \textbf{forwiernan}'s selected government.
% mdparagraph:P262:end

\end{minipage}\par\addvspace{4pt}

\par\begin{minipage}{\linewidth}\raggedright
% mdparagraph:P263:begin
\textbf{III.15.5 (EX-III-15-05).}
% mdparagraph:P263:end

% item: EX-III-15-05
% key: KEY-EX-III-15-05
% mdparagraph:P264:begin
\emph{The monk did it in Christ’s name.}
% mdparagraph:P264:end

% mdparagraph:P265:begin
\textbf{Se munuc dyde hit on Crīstes naman.}
% mdparagraph:P265:end

% mdparagraph:P266:begin
\textbf{Target and interpretation:} The monk performed the deed, represented by accusative \textbf{hit}. In the bounded name-expression, \textbf{naman} is dative after \textbf{on}, while \textbf{Crīstes} is its dependent genitive: the deed was done in Christ's name. The whole phrase's external case and its possessor's internal case remain distinct. Do not make both nouns genitive or replace \textbf{naman} by nominative \textbf{nama}.
% mdparagraph:P266:end

\end{minipage}\par\addvspace{4pt}

\par\begin{minipage}{\linewidth}\raggedright
% mdparagraph:P267:begin
\textbf{III.15.6 (EX-III-15-06).}
% mdparagraph:P267:end

% item: EX-III-15-06
% key: KEY-EX-III-15-06
% mdparagraph:P268:begin
\emph{He will separate the goats from the sheep.}
% mdparagraph:P268:end

% mdparagraph:P269:begin
\textbf{Hē tōtwǣmð ðā gǣt fram ðǣm scēapum.}
% mdparagraph:P269:end

% mdparagraph:P270:begin
\textbf{Target and interpretation:} Singular \textbf{hē} will separate accusative plural \textbf{ðā gǣt} from dative plural \textbf{ðǣm scēapum}. \textbf{Tōtwǣmð} is present third singular with contextual future reference. The unit explicitly assigns this separation frame for animals. \textbf{Gāt} would make the direct object singular; \textbf{scēap} after \textbf{fram} would omit the required dative plural ending.
% mdparagraph:P270:end

\end{minipage}\par\addvspace{4pt}

\par\begin{minipage}{\linewidth}\raggedright
% mdparagraph:P271:begin
\textbf{III.15.7 (EX-III-15-07).}
% mdparagraph:P271:end

% item: EX-III-15-07
% key: KEY-EX-III-15-07
% mdparagraph:P272:begin
\emph{He poured her out wine and water.}
% mdparagraph:P272:end

% mdparagraph:P273:begin
\textbf{Hē scęncte hiere wīn and wæter.}
% mdparagraph:P273:end

% mdparagraph:P274:begin
\textbf{Target and interpretation:} He poured or served wine and water for her. \textbf{Scęncte} is singular preterite; \textbf{hiere} is the dative recipient. \textbf{Wīn} and \textbf{wæter} are coordinated accusative mass objects, not the subjects of a plural verb. The pronoun is governed by the serving frame, not a possessive modifying the drinks.
% mdparagraph:P274:end

\end{minipage}\par\addvspace{4pt}

\par\begin{minipage}{\linewidth}\raggedright
% mdparagraph:P275:begin
\textbf{III.15.8 (EX-III-15-08).}
% mdparagraph:P275:end

% item: EX-III-15-08
% key: KEY-EX-III-15-08
% mdparagraph:P276:begin
\emph{The Lord will come with his angels.}
% mdparagraph:P276:end

% mdparagraph:P277:begin
\textbf{Se dryhten cymð mid his ęnglum.}
% mdparagraph:P277:end

% mdparagraph:P278:begin
\textbf{Target and interpretation:} The Lord is the singular subject of future-reference \textbf{cymð}. Accompanying angels are dative plural \textbf{ęnglum} after \textbf{mid}, using the syncopated stem. \textbf{His} is the invariant third-person possessor and refers to the singular Lord; it is not made plural to match the angels. Nominative \textbf{ęnglas} after \textbf{mid} would lose the accompaniment construction.
% mdparagraph:P278:end

\end{minipage}\par\addvspace{4pt}

\par\begin{minipage}{\linewidth}\raggedright
% mdparagraph:P279:begin
\textbf{III.15.9 (EX-III-15-09).}
% mdparagraph:P279:end

% item: EX-III-15-09
% key: KEY-EX-III-15-09
% mdparagraph:P280:begin
\emph{The morning cometh, and also the night.}
% mdparagraph:P280:end

% mdparagraph:P281:begin
\textbf{Se morgen cymð and ēac sēo niht.}
% mdparagraph:P281:end

% mdparagraph:P282:begin
\textbf{Target and interpretation:} Morning comes, and night also comes. \textbf{Morgen} is masculine and \textbf{niht} feminine, hence \textbf{se} and \textbf{sēo}. The second coordinated clause leaves its singular coming predicate understood; repeating \textbf{cymð} there is also valid. Do not force a plural verb into the first clause merely because the sentence mentions two arrivals, or treat archaic English \emph{cometh} as a command.
% mdparagraph:P282:end

\end{minipage}\par\addvspace{4pt}

\par\begin{minipage}{\linewidth}\raggedright
% mdparagraph:P283:begin
\textbf{III.15.10 (EX-III-15-10).}
% mdparagraph:P283:end

% item: EX-III-15-10
% key: KEY-EX-III-15-10
% mdparagraph:P284:begin
\emph{Work the works of God.}
% mdparagraph:P284:end

% mdparagraph:P285:begin
\textbf{Wyrceað Godes weorc.}
% mdparagraph:P285:end

% mdparagraph:P286:begin
\textbf{Target and interpretation:} \textbf{Godes} is genitive singular dependent on accusative plural \textbf{weorc}, the works to be performed. The unchanged neuter noun shape can be plural. The model uses plural command \textbf{wyrceað}; singular \textbf{Wyrc Godes weorc} is equally defensible because the English does not specify the addressee's number. The genitive identifies whose works they are; it does not make God the grammatical subject of the command.
% mdparagraph:P286:end

\end{minipage}\par\addvspace{4pt}

\par\begin{minipage}{\linewidth}\raggedright
% mdparagraph:P287:begin
\textbf{III.15.11 (EX-III-15-11).}
% mdparagraph:P287:end

% item: EX-III-15-11
% key: KEY-EX-III-15-11
% mdparagraph:P288:begin
\emph{Clothe (sing.) the naked.}
% mdparagraph:P288:end

% mdparagraph:P289:begin
\textbf{Scrȳd ðā nacodan.}
% mdparagraph:P289:end

% mdparagraph:P290:begin
\textbf{Target and interpretation:} Singular \textbf{scrȳd} addresses one person. Weak accusative plural \textbf{ðā nacodan} names the naked people to be clothed. The command's addressee and its object have separate numbers: plural people do not call for a plural imperative. A singular object article would change the model's collective reference.
% mdparagraph:P290:end

\end{minipage}\par\addvspace{4pt}

\par\begin{minipage}{\linewidth}\raggedright
% mdparagraph:P291:begin
\textbf{III.15.12 (EX-III-15-12).}
% mdparagraph:P291:end

% item: EX-III-15-12
% key: KEY-EX-III-15-12
% mdparagraph:P292:begin
\emph{I will not sit on the king’s throne.}
% mdparagraph:P292:end

% mdparagraph:P293:begin
\textbf{Ic ne sitte on ðæs cyninges setle.}
% mdparagraph:P293:end

% mdparagraph:P294:begin
\textbf{Target and interpretation:} First-singular \textbf{sitte} has future reference under negation. The throne is dative \textbf{setle} after locative \textbf{on}; \textbf{ðæs cyninges} is the internal genitive identifying its possessor. If a refusal or unwillingness is intended, \textbf{Ic ne wille on ðæs cyninges setle sittan} is also valid, using the taught plain infinitive. The present model is not mandatory for every reading of English \emph{will not}. Do not change the head noun to genitive because its possessor is genitive.
% mdparagraph:P294:end

\end{minipage}\par\addvspace{4pt}

\par\begin{minipage}{\linewidth}\raggedright
% mdparagraph:P295:begin
\textbf{III.15.13 (EX-III-15-13).}
% mdparagraph:P295:end

% item: EX-III-15-13
% key: KEY-EX-III-15-13
% mdparagraph:P296:begin
\emph{The bird sees the light.}
% mdparagraph:P296:end

% mdparagraph:P297:begin
\textbf{Se fugol gesiehð ðæt lēoht.}
% mdparagraph:P297:end

% mdparagraph:P298:begin
\textbf{Target and interpretation:} Nominative masculine \textbf{se fugol} is the perceiver; neuter accusative \textbf{ðæt lēoht} is the light seen. \textbf{Gesiehð} is present third singular. The object has the same surface as a nominative neuter phrase, but its role follows the clause. No extra finite or infinitival complement is needed when only the light is named.
% mdparagraph:P298:end

\end{minipage}\par\addvspace{4pt}

\par\begin{minipage}{\linewidth}\raggedright
% mdparagraph:P299:begin
\textbf{III.15.14 (EX-III-15-14).}
% mdparagraph:P299:end

% item: EX-III-15-14
% key: KEY-EX-III-15-14
% mdparagraph:P300:begin
\emph{The men will see the stars.}
% mdparagraph:P300:end

% mdparagraph:P301:begin
\textbf{Ðā męnn gesēoð ðā steorran.}
% mdparagraph:P301:end

% mdparagraph:P302:begin
\textbf{Target and interpretation:} \textbf{Ðā męnn} is the nominative plural subject, and \textbf{ðā steorran} the accusative plural stars. \textbf{Gesēoð} is a finite plural present with future reference, not an infinitive. Do not replace it with \textbf{gesēon} solely because the English has \emph{will}, or use the singular \textbf{gesiehð} with the plural men.
% mdparagraph:P302:end

\end{minipage}\par\addvspace{4pt}

\par\begin{minipage}{\linewidth}\raggedright
% mdparagraph:P303:begin
\textbf{III.15.15 (EX-III-15-15).}
% mdparagraph:P303:end

% item: EX-III-15-15
% key: KEY-EX-III-15-15
% mdparagraph:P304:begin
\emph{As long as you (plur.) fed the hungry, you did it to the Lord.}
% mdparagraph:P304:end

% mdparagraph:P305:begin
\textbf{Swā lange swā gē fēddon ðā hungrigan, gē dydon hit ðǣm dryhtne.}
% mdparagraph:P305:end

% mdparagraph:P306:begin
\textbf{Target and interpretation:} Both occurrences of \emph{you} are plural: \textbf{gē fēddon} and \textbf{gē dydon}. The hungry are weak accusative plural \textbf{ðā hungrigan}; \textbf{hit} is the deed done, and \textbf{ðǣm dryhtne} the dative person to whom it is reckoned. The local help supplies \textbf{swā lange swā}. Cook's note17 explains its biblical \emph{inasmuch as} force: the deeds to the hungry count as done to the Lord. Duration is not a mandatory interpretation. Do not exchange the recipient and the deed, use singular verbs, or treat the supplied connection as an unaided temporal-clause target. [4]
% mdparagraph:P306:end

\end{minipage}\par\addvspace{4pt}

\par\begin{minipage}{\linewidth}\raggedright
% mdparagraph:P307:begin
\textbf{III.15.16 (EX-III-15-16).}
% mdparagraph:P307:end

% item: EX-III-15-16
% key: KEY-EX-III-15-16
% mdparagraph:P308:begin
\emph{When did you see the stranger sick?}
% mdparagraph:P308:end

% mdparagraph:P309:begin
\textbf{Hwonne gesāwon gē ðone cuman untrumne?}
% mdparagraph:P309:end

% mdparagraph:P310:begin
\textbf{Target and interpretation:} The model uses plural \textbf{gē} with preterite \textbf{gesāwon}. The English leaves \emph{you} unmarked, so singular \textbf{Hwonne gesāwe ðū ðone cuman untrumne?} is also valid; a consistent dual interpretation is possible where two addressees are understood. The stranger and his perceived condition are both masculine accusative singular. Strong \textbf{untrumne} is an object predicate; a weak adjective inside an attributive \emph{the sick stranger} phrase would not demonstrate this stated construction.
% mdparagraph:P310:end

\end{minipage}\par\addvspace{4pt}

\par\begin{minipage}{\linewidth}\raggedright
% mdparagraph:P311:begin
\textbf{III.15.17 (EX-III-15-17).}
% mdparagraph:P311:end

% item: EX-III-15-17
% key: KEY-EX-III-15-17
% mdparagraph:P312:begin
\emph{The hungry man came to the ships.}
% mdparagraph:P312:end

% mdparagraph:P313:begin
\textbf{Se hungriga mann cōm tō ðǣm scipum.}
% mdparagraph:P313:end

% mdparagraph:P314:begin
\textbf{Target and interpretation:} The man is singular subject of past \textbf{cōm}. \textbf{Hungriga} is weak masculine nominative singular after \textbf{se}. The ships are dative plural \textbf{ðǣm scipum}, governed by destination \textbf{tō}, not objects of a transitive \emph{come}. Strong \textbf{hungrig} would miss the selected definite attributive pattern; accusative \textbf{scipu} would miss the destination frame.
% mdparagraph:P314:end

\end{minipage}\par\addvspace{4pt}

\par\begin{minipage}{\linewidth}\raggedright
% mdparagraph:P315:begin
\textbf{III.15.18 (EX-III-15-18).}
% mdparagraph:P315:end

% item: EX-III-15-18
% key: KEY-EX-III-15-18
% mdparagraph:P316:begin
\emph{The soldier brought us the horses’ fodder.}
% mdparagraph:P316:end

% mdparagraph:P317:begin
\textbf{Se cęmpa tōgelǣdde ūs ðāra horsa fōdor.}
% mdparagraph:P317:end

% mdparagraph:P318:begin
\textbf{Target and interpretation:} \textbf{Tōgelǣdde} is singular preterite agreeing with the soldier. Dative \textbf{ūs} names the recipients; accusative singular \textbf{fōdor} is what was brought. Inside that object phrase, genitive plural \textbf{ðāra horsa} identifies the horses whose fodder it is. The horses themselves are not the things brought. Keep the head's accusative separate from its dependent's genitive and from the recipients' dative.
% mdparagraph:P318:end

\end{minipage}\par\addvspace{4pt}

\par\begin{minipage}{\linewidth}\raggedright
% mdparagraph:P319:begin
\textbf{III.15.19 (EX-III-15-19).}
% mdparagraph:P319:end

% item: EX-III-15-19
% key: KEY-EX-III-15-19
% mdparagraph:P320:begin
\emph{The farmer has the blacksmith’s hammer.}
% mdparagraph:P320:end

% mdparagraph:P321:begin
\textbf{Se ierðling hæfð ðæs smiðes slęcge.}
% mdparagraph:P321:end

% mdparagraph:P322:begin
\textbf{Target and interpretation:} The farmer possesses accusative feminine singular \textbf{slęcge}. Its dependent \textbf{ðæs smiðes} is genitive masculine singular, identifying the blacksmith. The article agrees with the genitive possessor, not with the feminine hammer. Do not make the hammer genitive or the blacksmith an additional accusative object.
% mdparagraph:P322:end

\end{minipage}\par\addvspace{4pt}

\par\begin{minipage}{\linewidth}\raggedright
% mdparagraph:P323:begin
\textbf{III.15.20 (EX-III-15-20).}
% mdparagraph:P323:end

% item: EX-III-15-20
% key: KEY-EX-III-15-20
% mdparagraph:P324:begin
\emph{Who has the tailor’s needles?}
% mdparagraph:P324:end

% mdparagraph:P325:begin
\textbf{Hwā hæfð ðæs sēameres nǣdla?}
% mdparagraph:P325:end

% mdparagraph:P326:begin
\textbf{Target and interpretation:} Nominative \textbf{hwā} is the person who has the needles; \textbf{hæfð} is the finite singular verb. \textbf{Nǣdla} is accusative feminine plural, with genitive singular possessor \textbf{ðæs sēameres}. The original says \textbf{tailor's}, not sailor's. A singular needle or a different possessor loses the requested meaning; the question word is not accusative simply because something is possessed.
% mdparagraph:P326:end

\end{minipage}\par\addvspace{4pt}

\par\begin{minipage}{\linewidth}\raggedright
% mdparagraph:P327:begin
\textbf{III.15.21 (EX-III-15-21).}
% mdparagraph:P327:end

% item: EX-III-15-21
% key: KEY-EX-III-15-21
% mdparagraph:P328:begin
\emph{I hear the bellows (169) blow.}
% mdparagraph:P328:end

% mdparagraph:P329:begin
\textbf{Ic hīere ðā bęlgas blāwan.}
% mdparagraph:P329:end

% mdparagraph:P330:begin
\textbf{Target and interpretation:} \textbf{Ic} is the subject of first-singular \textbf{hīere}. \textbf{Ðā bęlgas} is the accusative plural perceived actor, with plain infinitive \textbf{blāwan} naming what the bellows do. Cook's \textbf{(169)} points to the infinitive construction. The new hearing frame is explicitly taught; it is not assumed merely from the earlier seeing frame. A plural finite \textbf{blāwað} would replace the infinitive and change the construction.
% mdparagraph:P330:end

\end{minipage}\par\addvspace{4pt}

\par\begin{minipage}{\linewidth}\raggedright
% mdparagraph:P331:begin
\textbf{III.15.22 (EX-III-15-22).}
% mdparagraph:P331:end

% item: EX-III-15-22
% key: KEY-EX-III-15-22
% mdparagraph:P332:begin
\emph{Do (om.) you (sing.) see the hammers smite?}
% mdparagraph:P332:end

% mdparagraph:P333:begin
\textbf{Gesiehst ðū ðā slęcga slēan?}
% mdparagraph:P333:end

% mdparagraph:P334:begin
\textbf{Target and interpretation:} Follow the explicit omission of English \emph{Do}: \textbf{gesiehst} itself is the finite second-singular question verb, agreeing with \textbf{ðū}. Feminine accusative plural \textbf{ðā slęcga} is the perceived actor of plain \textbf{slēan}. Here striking is used without naming the thing struck. Do not add a dummy \textbf{dōn}, make \textbf{slēan} agree with the hammers, or replace singular \textbf{ðū} by plural \emph{you}.
% mdparagraph:P334:end

\end{minipage}\par\addvspace{4pt}

\clearpage
% mdparagraph:P335:begin
\section*{R. The stranger and the monk}
% mdparagraph:P335:end

\reserve{9\baselineskip}
\par\begin{minipage}{\linewidth}\raggedright
% mdparagraph:P336:begin
\textbf{R01. Pronouns in the stranger’s account.}
% mdparagraph:P336:end

% item: NEW-U10-R01
% key: KEY-NEW-U10-R01
% mdtable:T08:begin
\par\begingroup
\fontsize{10.5}{13.2}\selectfont
\setlength{\LTleft}{0pt}\setlength{\LTright}{\fill}
\noindent
\begin{tabular}{@{}>{\raggedright\arraybackslash}p{37.65mm}>{\raggedright\arraybackslash}p{28.75mm}>{\raggedright\arraybackslash}p{51.98mm}@{}}
\toprule
\textbf{Occurrence} & \textbf{Referent} & \textbf{Case and governing construction} \\
\midrule
hine, after the arrival & the stranger & Accusative singular experiencer with hyngrede \\
him, after sealde & the stranger & Dative singular recipient with sealde \\
mē, in the stranger’s quotation & the stranger speaking & Dative singular person with forwiernst \\
ðē, in the monk’s quotation & the stranger addressed & Accusative singular person fed, with fēde \\
\bottomrule
\end{tabular}
\par\endgroup

% mdtable:T08:end
% mdparagraph:P337:begin
\textbf{Target and interpretation:} All four pronouns refer to the same stranger, but their case follows the current construction. Quoted \textbf{mē} and \textbf{ðē} change speech role, not real-world referent: the stranger speaks first, and the monk then addresses him. Do not identify the monk as \textbf{ðē}, or make an accusative experiencer into a nominative subject because the English translation uses \emph{he}.
% mdparagraph:P337:end

\end{minipage}\par\addvspace{4pt}

\par\begin{minipage}{\linewidth}\raggedright
% mdparagraph:P338:begin
\textbf{R02. The account’s genitives and weak noun.}
% mdparagraph:P338:end

% item: NEW-U10-R02
% key: KEY-NEW-U10-R02
% mdparagraph:P339:begin
\textbf{Ðāra munuca} is a genitive plural whole selected by nominative singular \textbf{ān}: one of the monks. \textbf{Mętes} is genitive singular food governed directly by \textbf{forwiernst}. These are different reasons for the same case.
% mdparagraph:P339:end

% mdparagraph:P340:begin
\textbf{Cuma} in the first sentence is nominative singular with \textbf{sum}, the arriving subject. Each later \textbf{se cuma} is likewise nominative, as the stranger speaks and then eats. \textbf{Ðone cuman} after \textbf{geseah} is accusative singular, the person seen.
% mdparagraph:P340:end

% mdparagraph:P341:begin
\textbf{Target and interpretation:} Assign case from the actual relation, not from the ending alone. The genitive group's plurality does not make \textbf{ān} plural, and the weak noun's \textbf{-an} does not itself mean plural. Selected genitive plural \textbf{munuca} must not be replaced by a spelling guessed from another headword variant.
% mdparagraph:P341:end

\end{minipage}\par\addvspace{4pt}

\par\begin{minipage}{\linewidth}\raggedright
% mdparagraph:P342:begin
\textbf{R03. Changing the means phrase.}
% mdparagraph:P342:end

% item: NEW-U10-R03
% key: KEY-NEW-U10-R03
% mdparagraph:P343:begin
\textbf{Gōde} is strong masculine singular instrumental, agreeing with \textbf{męte} in the monk's bare means phrase. With \textbf{mid}, the phrase is:
% mdparagraph:P343:end

% mdparagraph:P344:begin
\textbf{mid gōdum męte} — with good food.
% mdparagraph:P344:end

% mdparagraph:P345:begin
\textbf{Target and interpretation:} The means relation stays the same, but \textbf{mid} governs dative, so the adjective changes from distinct instrumental \textbf{gōde} to dative \textbf{gōdum}. The noun's surface stays \textbf{męte}. Rewriting the rest of the reply is unnecessary; leaving \textbf{gōde} after \textbf{mid} would fail this conversion.
% mdparagraph:P345:end

\end{minipage}\par\addvspace{4pt}

\par\begin{minipage}{\linewidth}\raggedright
% mdparagraph:P346:begin
\textbf{R04. The stranger and the monk: full translation.}
% mdparagraph:P346:end

% item: NEW-U10-R04
% key: KEY-NEW-U10-R04
% mdparagraph:P347:begin
A certain stranger came to the inn. He was hungry, and he was sick. One of the monks saw the stranger. The monk gave him food and water. The stranger said, “You do not deny me food.” The monk said, “I feed you with good food.” The stranger ate the food and drank the water.
% mdparagraph:P347:end

% mdparagraph:P348:begin
\textbf{Target and interpretation:} Preserve the full arrival, exchange and resolution. The stranger is hungry, receives the food, speaks the first quotation, and finally eats and drinks; the monk gives and feeds. \emph{Guest} for \emph{stranger}, \emph{lodging-house} for \emph{inn}, and natural English order are acceptable. Hunger does not mean that someone made him hungry, and \textbf{mętes} denotes the thing not denied, not a possessor. Speaker and addressee change between the quotations; \textbf{mē} in the first and \textbf{ðē} in the second nevertheless refer to the same stranger, as R01 explains.
% mdparagraph:P348:end

\end{minipage}\par\addvspace{4pt}

\clearpage
% mdparagraph:P349:begin
\section*{S. The Day of Judgment}
% mdparagraph:P349:end

\reserve{5\baselineskip}
\par\begin{minipage}{\linewidth}\raggedright
% mdparagraph:P350:begin
\textbf{S01. The Judgment tableau: participants and clauses.}
% mdparagraph:P350:end

% item: NEW-U10-S01
% key: KEY-NEW-U10-S01
% mdparagraph:P351:begin
The \textbf{Son of Man} comes, accompanied by all the angels. \textbf{All nations} are gathered before him; he separates them into two groups, with a shepherd's separation of sheep and goats as the comparison. He then places the sheep at his right and the goats at his left.
% mdparagraph:P351:end

% mdtable:T09:begin
\par\begingroup
\fontsize{10.5}{13.2}\selectfont
\setlength{\LTleft}{0pt}\setlength{\LTright}{\fill}
\noindent
\begin{tabular}{@{}>{\raggedright\arraybackslash}p{24mm}>{\raggedright\arraybackslash}p{97.19mm}@{}}
\toprule
\textbf{Finite clause} & \textbf{Boundary and interpretation} \\
\midrule
1. Coming & Witodlice Mannes Bearn cymð on his mægenðrymme. The attached and ealle ęnglas samod mid him tō ðǣm miclan dōme leaves a coming predicate understood. \\
2. Sitting & Ðonne sitt hē on ðǣm setle his mægenðrymnesse. \\
3. Gathering & and bēoð gegaderode ætforan him ealle ðēoda. The subject follows its plural passive. \\
4. Dividing & and hē tōscǣt hīe on twā. \\
5. Comparison & swā-swā scēaphierde tōscǣt scēap fram gātum. \\
6. Placement & Ðonne gelōgað hē ðā scēap on his swīðran hand, and ðā gǣt on his winstran. The coordinated final object/\allowbreak{}side phrase shares gelōgað. \\
\bottomrule
\end{tabular}
\par\endgroup

% mdtable:T09:end
% mdparagraph:P352:begin
\textbf{Target and interpretation:} There are six overt finite verbs: \textbf{cymð, sitt, bēoð, tōscǣt, tōscǣt, gelōgað}. \textbf{Gegaderode} is a participle, not a seventh finite verb; the angels have an understood predicate, not another written one. In the passive, \textbf{ealle ðēoda} is nominative feminine plural, the nations gathered; \textbf{ealle} is the supplied feminine \textbf{-e} variant. The source does not name the gathering agent. Natural masculine \textbf{hē/\allowbreak{}him} refers to the male person called by the grammatically neuter title \textbf{Bearn}. Do not turn the angels into objects of \textbf{mid} or treat the nations as those doing the gathering.
% mdparagraph:P352:end

\end{minipage}\par\addvspace{4pt}

\par\begin{minipage}{\linewidth}\raggedright
% mdparagraph:P353:begin
\textbf{S02. Separation and placement.}
% mdparagraph:P353:end

% item: NEW-U10-S02
% key: KEY-NEW-U10-S02
% mdparagraph:P354:begin
\textbf{Hīe} is accusative plural, referring to the assembled nations and governed by the first \textbf{tōscǣt}. In the placement clause, \textbf{ðā scēap} is neuter accusative plural and \textbf{ðā gǣt} feminine accusative plural; both are objects of \textbf{gelōgað}.
% mdparagraph:P354:end

% mdparagraph:P355:begin
In the shepherd comparison, \textbf{scēap} is again neuter accusative plural, whereas \textbf{fram} governs feminine dative plural \textbf{gātum}.
% mdparagraph:P355:end

% mdparagraph:P356:begin
\textbf{Target and interpretation:} Distinguish the entities separated or placed from the subject performing those actions. Endingless \textbf{scēap} need not be singular; \textbf{gǣt} is the selected umlauted plural in its object position. Do not call these placement objects nominative because the plural article has the same spelling in both cases, or make \textbf{gātum} accusative after \textbf{fram}.
% mdparagraph:P356:end

\end{minipage}\par\addvspace{4pt}

\par\begin{minipage}{\linewidth}\raggedright
% mdparagraph:P357:begin
\textbf{S03. Then and just as.}
% mdparagraph:P357:end

% item: NEW-U10-S03
% key: KEY-NEW-U10-S03
% mdparagraph:P358:begin
Both \textbf{ðonne} occurrences mean \textbf{then}. Each sequences a new sentence in the tableau: then he sits, and then he places the animals. \textbf{Swā-swā} means \textbf{just as} and introduces the finite shepherd-comparison clause.
% mdparagraph:P358:end

% mdparagraph:P359:begin
\textbf{Target and interpretation:} A sequencing adverb does not itself subordinate the following clause. The two instances of \textbf{ðonne} do not form a \emph{when … then} correlation, and neither means comparative \emph{than}. The supplied \textbf{swā-swā} clause compares two separations; it is not another event that happens after the division of the nations.
% mdparagraph:P359:end

\end{minipage}\par\addvspace{4pt}

\par\begin{minipage}{\linewidth}\raggedright
% mdparagraph:P360:begin
\textbf{S04. The Judgment tableau: full translation.}
% mdparagraph:P360:end

% item: NEW-U10-S04
% key: KEY-NEW-U10-S04
% mdparagraph:P361:begin
Now the Son of Man will come in his majesty, and all the angels together with him to the great judgment. Then he will sit on the throne of his majesty, and all nations will be gathered before him, and he will separate them into two groups, just as a shepherd separates sheep from goats. Then he will place the sheep on his right hand, and the goats on his left.
% mdparagraph:P361:end

% mdparagraph:P362:begin
\textbf{Target and interpretation:} Translate the complete three-sentence tableau, ending before the royal speech. A consistent present rendering is also defensible; the judgment context supports the model's future. \emph{Seat/\allowbreak{}throne}, \emph{glory/\allowbreak{}majesty}, \emph{as/\allowbreak{}just as}, and \emph{at his right/\allowbreak{}left side} can vary naturally. English may repeat the understood coming of the angels and the final \emph{hand}. Preserve the passive gathering, the division into two groups, and the general shepherd comparison. The two nouns for majesty are distinct source words, and \textbf{setle his mægenðrymnesse} has the locally supported characterizing genitive. The right/\allowbreak{}left expressions use the supplied bounded side-frame; do not infer their case solely from motion or from weak \textbf{-an}. [5]
% mdparagraph:P362:end

\end{minipage}\par\addvspace{4pt}

\clearpage
% mdparagraph:P363:begin
\section*{Source and form notes}
% mdparagraph:P363:end

\begingroup\fontsize{10}{12.8}\selectfont
% mdparagraph:P364:begin
[1] Cook, \emph{Exercises in Old English} (1895), Group III, Exercises13–15, printed pp6–7 /\allowbreak{} PDF14–15. Punctuation spacing and typographic apostrophes follow the chapter's declared transcription policy. All 55 numbered prompts are retained.
% mdparagraph:P364:end

% mdparagraph:P365:begin
[2] Cook, \emph{A First Book in Old English}, 2nd ed. (1895), §43.8, printed p30, for \textbf{fæder}, and §52, printed p36, for \textbf{gāt}. Selected forms and admitted variants are distinguished in the tables; neither table is an unrestricted historical reconstruction.
% mdparagraph:P365:end

% mdparagraph:P366:begin
[3] Cook, §129, printed p77, and §193a, printed p103; Fulk, \emph{An Introductory Grammar of Old English} (2020), §123, printed pp89–90. The supported \textbf{unne} is common to the cited treatments; their imperative/\allowbreak{}subjunctive terminology differs. The indicative second-singular gap is limited to the evidence used for this course.
% mdparagraph:P366:end

% mdparagraph:P367:begin
[4] Cook, §§174–178, printed pp97–99, for instrumental usage. The right-hand expression is corroborated by \textbf{standað} in the unassigned continuation on p134; its static verb does not change the attested accusative. The local drinking frame comes from Reader III, p136, editorial lines5–6. The \textbf{swā lange swā} help retains p135, original note17, including the \emph{inasmuch as} interpretation. These corroborating spans do not enlarge the assigned reading.
% mdparagraph:P367:end

% mdparagraph:P368:begin
[5] Cook, Reader III, \emph{The Day of Judgment}, printed p134 /\allowbreak{} PDF150, the selected opening three sentences, editorial prose lines1–7 anchored by the printed marginal5. The exact source text and its fourteen original notes remain separate from the nine local help groups and from this new English translation.
% mdparagraph:P368:end

\endgroup
\end{document}
